% ============================================================
% TARDUS Protocol Specification
% Phase 0 — Version 0.1
% ============================================================
\documentclass[11pt,letterpaper]{article}

% --- packages ---
\usepackage[T1]{fontenc}
\usepackage[utf8]{inputenc}
\usepackage[english]{babel}
\usepackage{amsmath, amsthm, amssymb, amsfonts}
\usepackage{mathtools}
\usepackage{xcolor}
\usepackage{microtype}
\usepackage[margin=1in]{geometry}
\usepackage{hyperref}
\usepackage{cleveref}
\usepackage{enumitem}
\usepackage{booktabs}
\usepackage{longtable}
\usepackage{listings}
\usepackage{tikz}
\usepackage{fancyhdr}

\hypersetup{
  colorlinks=true,
  linkcolor=blue!50!black,
  citecolor=green!50!black,
  urlcolor=blue!70!black,
  pdftitle={TARDUS Protocol Specification},
  pdfauthor={TARDUS Architecture Committee},
  pdfsubject={Fully private payment protocol on the Solana main chain},
  pdfkeywords={Solana, blind signature, threshold, privacy, Chaum, refresh}
}

% --- bibliography ---
\bibliographystyle{plain}

% --- theorem environments ---
\newtheorem{theorem}{Theorem}[section]
\newtheorem{lemma}[theorem]{Lemma}
\newtheorem{proposition}[theorem]{Proposition}
\newtheorem{corollary}[theorem]{Corollary}

\theoremstyle{definition}
\newtheorem{definition}[theorem]{Definition}
\newtheorem{example}[theorem]{Example}
\newtheorem{construction}[theorem]{Construction}

\theoremstyle{remark}
\newtheorem{remark}[theorem]{Remark}
\newtheorem*{notation}{Notation}

% cleveref name aliases (so \Cref recognises our theorem-like environments)
\crefname{theorem}{Theorem}{Theorems}
\crefname{lemma}{Lemma}{Lemmas}
\crefname{proposition}{Proposition}{Propositions}
\crefname{corollary}{Corollary}{Corollaries}
\crefname{definition}{Definition}{Definitions}
\crefname{construction}{Construction}{Constructions}
\crefname{example}{Example}{Examples}
\crefname{remark}{Remark}{Remarks}

% --- notation macros ---
\newcommand{\negl}{\mathsf{negl}}
\newcommand{\poly}{\mathsf{poly}}
\newcommand{\ROM}{\mathsf{ROM}}
\newcommand{\Adv}{\mathcal{A}}
\newcommand{\Sim}{\mathcal{S}}
\newcommand{\AdvG}[1]{\mathsf{Game}_{#1}}

% TARDUS-specific
\newcommand{\Mint}{\mathsf{Mint}}
\newcommand{\Coin}{\mathsf{Coin}}
\newcommand{\Refresh}{\mathsf{Refresh}}
\newcommand{\Spend}{\mathsf{Spend}}
\newcommand{\Vault}{\mathsf{Vault}}
\newcommand{\Commit}{\mathsf{Commit}}
\newcommand{\Nullifier}{\mathsf{Nullifier}}
\newcommand{\Denom}{\mathsf{Denom}}
\newcommand{\Keyset}{\mathsf{Keyset}}
\newcommand{\Epoch}{\mathsf{Epoch}}

% --- title ---
\title{
  \vspace{-2em}
  \Huge\textbf{TARDUS} \\[0.8em]
  \Large A Fully Private Payment Protocol \\[0.3em] on Solana --- Specification \\[1.2em]
  \normalsize Phase 0 --- Version 0.1 (draft)
}
\author{
  TARDUS Architecture Committee \\
  \texttt{nzengi}
}
\date{2026-05-20}

% --- page style ---
\pagestyle{fancy}
\fancyhf{}
\rhead{\small TARDUS Spec v0.1}
\lhead{\small\leftmark}
\rfoot{\small\thepage}

\begin{document}

\maketitle
\thispagestyle{empty}

\begin{abstract}
\noindent
TARDUS is a fully private bearer payment protocol on the Solana main chain. It hides the sender, the receiver, the amount, and the asset type. It uses neither zero-knowledge proofs, nor pairing-based cryptography, nor per-circuit trusted setup. Instead, it builds on David Chaum's 1982 blind-signature bearer cash, fused with the operational hygiene of the 2020s: threshold blind Schnorr signatures (\cite{stinson-strobl2001} and \cite{komlo-goldberg2020}), a $\kappa$-fold cut-and-choose refresh protocol (\cite{dold2019}), a mint committee with mandatory FIPS 140-2 Level 3 HSMs, proactive secret-sharing rotation (\cite{gennaro-jarecki-krawczyk-rabin1999}), and a transparent collateral vault on top of Solana Token-2022. Computational unforgeability follows from the Pointcheval-Stern (\cite{pointcheval-stern1996}) random-oracle reduction; statistical blindness holds at issuance; double-spend resistance follows from an on-chain mathematical collateral invariant. This document is the Phase 0 specification; it serves as the foundation for the Phase 1 reference implementation, Phase 2-3 operations and wallet, Phase 4 audit and formal verification, and Phase 5 mainnet deployment.
\end{abstract}

\vspace{1em}

\begin{center}
\fbox{\parbox{0.9\textwidth}{\small
\textbf{License:} This spec and all accompanying artifacts are released under TARDUS-PROPRIETARY-1.0. Visibility is for reference only; copying, derivation, redistribution, and AI training/inference are prohibited. See the \texttt{LICENSE} file at the repository root.
}}
\end{center}

\vspace{1em}

\tableofcontents

\newpage

% --- sections ---
\section{Introduction}
\label{sec:introduction}

\textit{This section will be authored as a Phase 0 sub-task in \texttt{spec/sections/01-introduction.tex}. Scope: problem statement, threat model, TARDUS's contribution, comparison with adjacent systems (Cashu, Taler, Zcash, Aztec), document organization.}

\section{Cryptographic Primitives}
\label{sec:primitives}

This section establishes the cryptographic primitives, notation, and security assumptions underlying TARDUS. All subsequent protocols (\Cref{sec:mint,sec:refresh,sec:onchain}) reference the definitions and constructions introduced here.

\subsection{Notation}

We write $\lambda \in \mathbb{N}$ for the security parameter (default $\lambda = 128$ for TARDUS v1). A function $f \colon \mathbb{N} \to [0,1]$ is \emph{negligible}, written $f(\lambda) = \negl(\lambda)$, if for every polynomial $p$ there exists $\lambda_0$ such that for all $\lambda \geq \lambda_0$, $f(\lambda) < 1/p(\lambda)$.

Probabilistic polynomial-time is abbreviated \emph{PPT}. We write $x \leftarrow_\$ X$ to denote sampling $x$ uniformly at random from a set $X$, and $y \leftarrow A(x)$ to denote running an algorithm $A$ on input $x$ and assigning the output to $y$. We use $[n] = \{1, 2, \ldots, n\}$.

\subsection{Elliptic-Curve Group}
\label{sec:group}

TARDUS operates over the prime-order subgroup of the Edwards curve \texttt{edwards25519}, the group used by Ed25519. We denote:

\begin{itemize}[leftmargin=2em,itemsep=0pt]
  \item $\mathbb{G}$: the prime-order subgroup, of order $\ell = 2^{252} + 27742317777372353535851937790883648493$.
  \item $G \in \mathbb{G}$: the standard base point of \texttt{edwards25519}.
  \item $\mathbb{F}_\ell$: the scalar field, integers modulo $\ell$.
\end{itemize}

We assume the \emph{Elliptic-Curve Discrete Logarithm Problem (ECDLP)} is hard in $\mathbb{G}$: for any PPT adversary $\Adv$,
\[
  \Pr\!\bigl[\, x \leftarrow_\$ \mathbb{F}_\ell;\; X \gets xG;\; x' \leftarrow \Adv(\mathbb{G}, G, X) \;:\; x' = x \,\bigr] = \negl(\lambda).
\]
This assumption underlies every signature unforgeability claim in TARDUS.

\subsection{Hash Functions and the Random Oracle Model}
\label{sec:hash}

TARDUS uses two concrete hash functions, each with a distinct purpose:
\begin{itemize}[leftmargin=2em,itemsep=0pt]
  \item SHA-256~\cite{nist-sp800-185} for keyset identifiers, commitment binding, and any context whose output is consumed as bytes (not as a scalar).
  \item SHA-512~\cite{nist-sp800-185} for the hash-to-scalar map $H_{\mathbb{F}_\ell}$ used in Schnorr challenges (\Cref{con:schnorr,con:blind-schnorr,con:threshold-schnorr,con:blind-threshold}).
\end{itemize}

For security analysis we model both as \emph{random oracles} (the ROM, following~\cite{pointcheval-stern1996}).

The hash-to-scalar map $H_{\mathbb{F}_\ell} \colon \{0,1\}^* \to \mathbb{F}_\ell$ is realised by SHA-512 reduced modulo $\ell$. The reduction bias is bounded by $\ell / 2^{512} \approx 2^{252} / 2^{512} = 2^{-260}$, which is negligible. A naive 256-bit hash reduced modulo $\ell$ would carry a bias of $\approx 2^{-4} \approx 6\%$ and is therefore \emph{not} used for scalar derivation anywhere in the protocol.

\subsection{Schnorr Signatures}
\label{sec:schnorr}

The Schnorr signature scheme~\cite{schnorr1991} over $\mathbb{G}$ is parameterised by $(\mathbb{G}, G, \ell, H_{\mathbb{F}_\ell})$.

\begin{construction}[Schnorr signature]
\label{con:schnorr}
\begin{description}[leftmargin=1.5em,style=nextline,itemsep=0.2em]
  \item[$\mathsf{KeyGen}(1^\lambda)$:] Sample $sk \leftarrow_\$ \mathbb{F}_\ell$; compute $pk \gets sk \cdot G$; return $(pk, sk)$.
  \item[$\mathsf{Sign}(sk, m)$:] Sample $k \leftarrow_\$ \mathbb{F}_\ell$; compute $R \gets kG$;
    \[ c \gets H_{\mathbb{F}_\ell}(R \,\|\, pk \,\|\, m), \quad s \gets k + c \cdot sk \pmod{\ell}. \]
    Return $\sigma = (R, s)$.
  \item[$\mathsf{Verify}(pk, m, \sigma = (R, s))$:] Compute $c \gets H_{\mathbb{F}_\ell}(R \,\|\, pk \,\|\, m)$; return $\mathbf{1}[\,sG = R + c \cdot pk\,]$.
\end{description}
\end{construction}

\begin{theorem}[Pointcheval-Stern, informal]
\label{thm:schnorr-euf}
Schnorr signatures are existentially unforgeable under chosen-message attack in the random-oracle model, by reduction to ECDLP via the forking lemma~\cite{pointcheval-stern1996}.
\end{theorem}

\subsection{Blind Schnorr Signatures}
\label{sec:blind-schnorr}

The blind variant allows a signer to issue a signature on a message $m$ without learning $m$, the resulting signature, or the link between issuance and unblinding.

\begin{construction}[Blind Schnorr signature]
\label{con:blind-schnorr}
The blind issuance protocol between user $U$ (who knows $m$) and signer $S$ (who holds $sk$ with public key $pk = sk \cdot G$):
\begin{enumerate}[leftmargin=2em,itemsep=0.2em]
  \item $S$: sample $k \leftarrow_\$ \mathbb{F}_\ell$; compute $R \gets kG$; send $R$ to $U$.
  \item $U$: sample $\alpha, \beta \leftarrow_\$ \mathbb{F}_\ell$; compute
    \[
      R' \gets R + \alpha G + \beta \cdot pk, \quad
      c' \gets H_{\mathbb{F}_\ell}(R' \,\|\, pk \,\|\, m), \quad
      c \gets c' + \beta \pmod{\ell}.
    \]
    Send $c$ to $S$.
  \item $S$: compute $s \gets k + c \cdot sk \pmod{\ell}$; send $s$ to $U$.
  \item $U$: compute $s' \gets s + \alpha \pmod{\ell}$; output $\sigma' = (R', s')$ as the unblinded signature on $m$.
\end{enumerate}
\end{construction}

A direct calculation shows $s' G = R' + c' \cdot pk$, so $\sigma'$ is a valid Schnorr signature on $m$ under $pk$ in the sense of~\Cref{con:schnorr}.

\begin{theorem}[Statistical blindness]
\label{thm:blindness}
For every (possibly malicious) PPT signer $S^*$, the joint distribution of issuance transcripts $(R, c, s)$ observed by $S^*$ is statistically independent of the unblinded signature $(R', s', m)$. The unblinding factors $(\alpha, \beta)$ are uniform in $\mathbb{F}_\ell^2$, inducing a uniform marginal over all possible $m$ consistent with the transcript.
\end{theorem}

\begin{theorem}[One-more unforgeability~\cite{pointcheval-stern1996}]
\label{thm:bs-unforgeability}
If $\Adv$ makes $q$ valid blind-issuance queries, the probability that $\Adv$ outputs $q+1$ valid distinct $(m, \sigma')$ pairs is bounded by $q \cdot \mathsf{Adv}^{\mathsf{ECDLP}}_{\mathcal{B}}(\lambda) + \negl(\lambda)$ for some PPT reduction $\mathcal{B}$ in the ROM.
\end{theorem}

\subsection{Pedersen Verifiable Secret Sharing}
\label{sec:vss}

For threshold key generation, TARDUS uses Pedersen's non-interactive VSS~\cite{pedersen1991}, which provides information-theoretic hiding of the secret while allowing public verification of share consistency.

\begin{construction}[Pedersen VSS]
\label{con:vss}
Given $\mathbb{G}$ with generator $G$ and an independent generator $H$ (with $\log_G H$ unknown), a dealer with secret $s \in \mathbb{F}_\ell$ shares $s$ among $n$ parties with threshold $t$:
\begin{enumerate}[leftmargin=2em,itemsep=0.2em]
  \item Sample $f(x) = s + a_1 x + \cdots + a_{t-1} x^{t-1}$, where $a_i \leftarrow_\$ \mathbb{F}_\ell$.
  \item Sample $g(x) = r + b_1 x + \cdots + b_{t-1} x^{t-1}$, where $r, b_i \leftarrow_\$ \mathbb{F}_\ell$.
  \item Publish commitments $C_i = a_i G + b_i H$ for $i = 0, \ldots, t-1$ (with $a_0 = s, b_0 = r$).
  \item Send share $(s_j, r_j) = (f(j), g(j))$ privately to party $j$.
\end{enumerate}
Party $j$ verifies by checking $s_j G + r_j H = \sum_{i=0}^{t-1} j^i C_i$.
\end{construction}

This scheme provides \emph{binding} (the dealer cannot equivocate on $s$ after publishing the $C_i$) under the discrete-log assumption, and \emph{statistical hiding} (any set of fewer than $t$ shares reveals no information about $s$).

\subsection{Distributed Key Generation}
\label{sec:dkg}

TARDUS bootstraps the threshold mint key via a DKG ceremony following Gennaro et al.~\cite{gennaro-jarecki-krawczyk-rabin1999} with the round-optimisation of Komlo-Goldberg~\cite{komlo-goldberg2020}.

\begin{construction}[FROST-style DKG, simplified]
\label{con:dkg}
For $n$ participants with threshold $t$:
\begin{enumerate}[leftmargin=2em,itemsep=0.2em]
  \item Each participant $P_i$ acts as a Pedersen-VSS dealer for a freshly sampled secret $s_i \leftarrow_\$ \mathbb{F}_\ell$.
  \item Each $P_i$ broadcasts a Schnorr proof of knowledge of $s_i$ against the commitment $C_0^{(i)} = s_i G$.
  \item All commitments and proofs are gossiped; participants verify each other's commitments and proofs.
  \item Each $P_i$ aggregates the received shares: $sk_i = \sum_{j=1}^{n} s_{j \to i}$.
  \item The joint public key is $PK = \sum_{j=1}^{n} C_0^{(j)} = \bigl( \sum_{j=1}^{n} s_j \bigr) G$.
\end{enumerate}
The joint secret key is $SK = \sum_{j=1}^{n} s_j$, distributed across the $n$ participants as Shamir shares with threshold $t$.
\end{construction}

TARDUS adopts $n = 30$, $t = 14$ as the canonical committee size (\Cref{sec:mint}).

\subsection{Threshold Schnorr Signatures (Stinson-Strobl)}
\label{sec:threshold-schnorr}

Given a DKG-shared key, $t$ of $n$ participants can jointly produce a Schnorr signature on $m$ without ever reconstructing $SK$.

\begin{construction}[Threshold Schnorr~\cite{stinson-strobl2001}, simplified]
\label{con:threshold-schnorr}
For a signing set $S \subseteq [n]$ of size $|S| = t$:
\begin{enumerate}[leftmargin=2em,itemsep=0.2em]
  \item Each $P_i$ for $i \in S$ samples $k_i \leftarrow_\$ \mathbb{F}_\ell$, computes $R_i \gets k_i G$, broadcasts $R_i$.
  \item Aggregate: $R \gets \sum_{i \in S} R_i$.
  \item Challenge: $c \gets H_{\mathbb{F}_\ell}(R \,\|\, PK \,\|\, m)$.
  \item Each $P_i$ computes its partial signature $s_i \gets k_i + c \cdot \lambda_i^{(S)} \cdot sk_i \pmod{\ell}$, where $\lambda_i^{(S)}$ is the Lagrange coefficient of $i$ relative to $S$.
  \item Aggregate: $s \gets \sum_{i \in S} s_i$.
  \item Output $\sigma = (R, s)$, verifiable as a standard Schnorr signature under $PK$ via~\Cref{con:schnorr}.
\end{enumerate}
\end{construction}

\begin{theorem}[\cite{stinson-strobl2001}, informal]
\label{thm:threshold-euf}
The threshold Schnorr scheme of~\Cref{con:threshold-schnorr} is unforgeable against a static adversary corrupting fewer than $t$ participants, by reduction to ECDLP in the ROM, provided the underlying VSS is secure.
\end{theorem}

\subsection{Blind Threshold Schnorr --- the TARDUS Construction}
\label{sec:blind-threshold}

TARDUS composes blind issuance (\Cref{con:blind-schnorr}) with threshold signing (\Cref{con:threshold-schnorr}). The user blinds against the joint public key $PK$; the $t$-of-$n$ committee collaboratively produces a blind threshold signature; the user unblinds.

\begin{construction}[Blind threshold Schnorr]
\label{con:blind-threshold}
For a committee with joint key $PK$, threshold $t$, and signing set $S$ of size $t$:
\begin{enumerate}[leftmargin=2em,itemsep=0.2em]
  \item Each $P_i$ samples $k_i \leftarrow_\$ \mathbb{F}_\ell$ and broadcasts $R_i \gets k_i G$. The committee aggregates $R \gets \sum_{i \in S} R_i$ and sends $R$ to the user.
  \item User: sample $\alpha, \beta \leftarrow_\$ \mathbb{F}_\ell$; compute
    \[
      R' \gets R + \alpha G + \beta \cdot PK, \quad
      c' \gets H_{\mathbb{F}_\ell}(R' \,\|\, PK \,\|\, m), \quad
      c \gets c' + \beta \pmod{\ell}.
    \]
    Send $c$ to the committee.
  \item Each $P_i$ computes $s_i \gets k_i + c \cdot \lambda_i^{(S)} \cdot sk_i \pmod{\ell}$ and broadcasts.
  \item Aggregate: $s \gets \sum_{i \in S} s_i$; send $s$ to the user.
  \item User: $s' \gets s + \alpha \pmod{\ell}$; output $\sigma' = (R', s')$.
\end{enumerate}
\end{construction}

Correctness follows by combining the linearity arguments of~\Cref{con:blind-schnorr} and~\Cref{con:threshold-schnorr}. Security is the subject of~\Cref{sec:security} (theorems T1 and T2); the formal proofs are placeholders in \texttt{proofs/unforgeability.ec} and \texttt{proofs/blindness.ec} pending mechanisation in Phase~4.

\subsection{Solana Cryptographic Syscall Surface}
\label{sec:syscalls}

The on-chain program of~\Cref{sec:onchain} verifies signatures using the Solana \texttt{ed25519} signature verification path, which costs approximately 5\,000 compute units per signature verification. Hashing is performed via \texttt{sol\_sha256} (approximately 85\,CU per call) and, where state-tree compatibility requires it, \texttt{sol\_poseidon} (Light Protocol's BN254 Poseidon). TARDUS deliberately avoids the \texttt{alt\_bn128} pairing syscall, as no pairing-based primitive is used by the protocol; the only place a pairing-using algorithm appears anywhere in the stack is inside Light Protocol's validity-proof verifier for compressed accounts, which TARDUS depends on transitively (\Cref{sec:onchain}).

A per-instruction CU budget table appears in~\Cref{sec:onchain}.

\section{Threshold Mint Protocol}
\label{sec:mint}

The TARDUS mint is a $t$-of-$n$ threshold service that issues blind Schnorr signatures (\Cref{con:blind-threshold}) under a jointly-held key. No single operator can either forge a signature, learn the message being issued, or unilaterally revoke a denomination. This section specifies the committee topology, the distributed key generation ceremony, the threshold blind signing protocol, denomination key rotation, and the revoke/recoup paths.

\subsection{Committee Topology}
\label{sec:mint-topology}

The canonical TARDUS v1 committee is fixed at $n = 30$ validators with threshold $t = 14$. This choice reflects three operational constraints derived from production lessons (see \texttt{research/PRODUCTION\_LESSONS.md}):

\begin{itemize}[leftmargin=2em,itemsep=0pt]
  \item \textbf{Liveness margin.} With $n - t + 1 = 17$ validators required online for signing, the system tolerates up to $13$ simultaneous offline validators. At an operational availability of $99\%$ per validator, the probability that more than $13$ are simultaneously offline is $< 2^{-30}$, well below any realistic outage scenario.
  \item \textbf{Forgery resistance.} Forging a signature requires corrupting at least $14$ validators. Under independent compromise probability $p$ per validator, the joint compromise probability is bounded by $\binom{30}{14} p^{14} \approx 1.45 \times 10^8 \cdot p^{14}$; for $p \leq 10^{-3}$ this gives $\leq 2^{-12}$ per epoch.
  \item \textbf{Communication cost.} Each DKG and signing round produces $\mathcal{O}(n^2)$ point-to-point messages in the pessimistic case. At $n = 30$ this is $900$ messages per round, well within practical operational bounds.
\end{itemize}

\subsection{Three-Tier Key Architecture}
\label{sec:mint-keyarch}

Each validator $V_i$ holds three distinct cryptographic identities, separated by hardware and process boundary:

\begin{description}[leftmargin=1.5em,style=nextline,itemsep=0.3em]
  \item[Tier 1 --- Offline Master Identity Key (MIK)] A long-term Ed25519 keypair, generated once on an air-gapped device and stored exclusively in an offline HSM. The MIK signs (a) the validator's enrolment certificate, (b) all ceremony transcripts, and (c) the validator's Tier 2 keys with limited validity. The MIK never participates in online operations.
  \item[Tier 2 --- Online Threshold Helper (OTH)] An Ed25519 keypair issued by Tier 1, with a validity period of at most $90$ epochs. The OTH holds the validator's share $sk_i$ of the joint mint key inside an online FIPS 140-2 Level 3 HSM. All threshold signing operations are gated through the OTH and physically executed inside the HSM.
  \item[Tier 3 --- Public API Endpoint (PAE)] A network-facing daemon that receives user requests (blind issuance, refresh, recoup) and forwards them to Tier 2 via authenticated UNIX domain sockets. The PAE holds no secret key material; compromise of the PAE limits the attacker to denial of service for the duration of control.
\end{description}

The architecture mirrors the GNU Taler exchange's three-tier separation (\cite{taler-exchange-manual}) but elevates Tier 2 with mandatory HSM custody, addressing the gap noted in the Taler manual itself (``does not yet support the use of a hardware security module'').

\subsection{Hardware Security Module Requirements}
\label{sec:mint-hsm}

Each validator's Tier 2 HSM must satisfy the following requirements:

\begin{itemize}[leftmargin=2em,itemsep=0pt]
  \item Certified to FIPS 140-2 Level 3~\cite{fips140-2} or equivalent (Common Criteria EAL 4+).
  \item Native support for Ed25519 / edwards25519 scalar multiplication and signature operations.
  \item Capacity for at least $200$ key handles (to accommodate $LOOKAHEAD = 30$ epochs of pre-generated share material with rotation overlap).
  \item Tamper-evident enclosure with audit logging of all key operations.
  \item Cryptographically authenticated firmware (vendor-signed updates only).
\end{itemize}

Reference implementations targeting YubiHSM 2, Thales Luna T-Series, and AWS CloudHSM are produced in Phase 2 (see \texttt{research/PRODUCTION\_LESSONS.md} §L7).

\subsection{Distributed Key Generation Ceremony}
\label{sec:mint-dkg}

The DKG ceremony bootstraps the joint mint key. It follows~\cite{gennaro-jarecki-krawczyk-rabin1999} with the round optimisation of~\cite{komlo-goldberg2020}, instantiated over \texttt{edwards25519} using the primitives of~\Cref{sec:primitives}.

\subsubsection{Ceremony Identifier}
\label{sec:mint-ceremony-id}

Each ceremony is parameterised by a $128$-bit \emph{ceremony identifier}
\[
  \mathsf{ceremony\_id} \in \{0,1\}^{128},
\]
generated as $\mathsf{ceremony\_id} = \mathsf{HKDF}(\mathsf{epoch} \,\|\, \mathsf{nonce})$ where $\mathsf{nonce}$ is $64$ bits of fresh randomness sampled jointly via verifiable random function output during ceremony initialisation. The identifier is a domain separator that binds every signed message in the ceremony to a unique session, preventing cross-ceremony replay.

\subsubsection{Round 1 --- Pedersen VSS Share Distribution}
\label{sec:mint-dkg-r1}

Each validator $V_i$, for $i \in [n]$, acts as a Pedersen-VSS dealer per~\Cref{con:vss}:

\begin{enumerate}[leftmargin=2em,itemsep=0.2em]
  \item $V_i$ samples a fresh polynomial $f_i(x) = s_i + a_{i,1} x + \cdots + a_{i,t-1} x^{t-1}$ over $\mathbb{F}_\ell$, with $s_i$ being the validator's contribution to the joint secret.
  \item $V_i$ samples a fresh blinding polynomial $g_i(x) = r_i + b_{i,1} x + \cdots + b_{i,t-1} x^{t-1}$.
  \item $V_i$ publishes both Pedersen commitments $C_{i,k} = a_{i,k} G + b_{i,k} H$ and Feldman commitments $A_{i,k} = a_{i,k} G$ for $k = 0, \ldots, t-1$ (with $a_{i,0} = s_i$, $b_{i,0} = r_i$). Here $H$ is an independent generator with unknown discrete logarithm relative to $G$, derived as $H = \mathsf{hash\_to\_point}(\text{``TARDUS-pedersen-H-v1''})$. The Pedersen set is used for share verification (statistical hiding of intermediate state); the Feldman set is used for joint-key derivation (\Cref{sec:mint-dkg-r3}) and as the public key against which the proof of knowledge is verified.
  \item $V_i$ broadcasts a Schnorr proof of knowledge $\pi_i$ of $s_i$ against the Feldman secret commitment $A_{i,0} = s_i G$.
  \item $V_i$ privately sends to each validator $V_j$ (for $j \neq i$) the shares $\sigma_{i \to j} = (f_i(j), g_i(j))$ over an authenticated, encrypted channel (Tier 1 MIK key-derived).
\end{enumerate}

Each validator $V_j$ verifies received shares by checking
\[
  f_i(j) \cdot G + g_i(j) \cdot H \stackrel{?}{=} \sum_{k=0}^{t-1} j^k \cdot C_{i,k}
\]
and verifies $\pi_i$ as a standard Schnorr signature (\Cref{con:schnorr}) over the implicit message $\mathsf{ceremony\_id} \,\|\, i$. Any failed verification opens a complaint in Round 2.

\subsubsection{Round 2 --- Complaint Phase}
\label{sec:mint-dkg-r2}

Validators with verification failures broadcast complaints of the form $(\mathsf{ceremony\_id}, V_j, V_i, \mathsf{evidence})$ where $V_j$ accuses $V_i$ of providing an invalid share. The accused validator $V_i$ must respond within the round 2 window by either:

\begin{itemize}[leftmargin=2em,itemsep=0pt]
  \item Revealing the share $\sigma_{i \to j}$ publicly (which the complainant or any party can verify against $V_i$'s commitments), or
  \item Failing to respond, in which case $V_i$ is removed from the qualified set $\mathsf{QUAL}$.
\end{itemize}

Public reveal of a share does not compromise security, because the underlying VSS commitments remain statistically hiding for any set of $< t$ revealed shares. The qualified set $\mathsf{QUAL} \subseteq [n]$ contains exactly the validators that survive Round 2.

\subsubsection{Round 3 --- Finalisation}
\label{sec:mint-dkg-r3}

The joint public key is computed deterministically from $\mathsf{QUAL}$ via the Feldman commitments:
\[
  PK = \sum_{i \in \mathsf{QUAL}} A_{i, 0} = \Bigl( \sum_{i \in \mathsf{QUAL}} s_i \Bigr) G.
\]
The Pedersen commitments $C_{i,0} = s_i G + r_i H$ are not used here, because their additive blinding component $\sum_i r_i H$ would obscure the joint key.

Each validator $V_j$ derives its share of the joint secret as
\[
  sk_j = \sum_{i \in \mathsf{QUAL}} f_i(j),
\]
which is stored inside the validator's Tier 2 HSM. The validator's auxiliary blinding share $\sum_{i \in \mathsf{QUAL}} g_i(j)$ is similarly retained.

For the ceremony to succeed, $|\mathsf{QUAL}| \geq t$ is required; otherwise the ceremony aborts and a new ceremony is initiated with a fresh $\mathsf{ceremony\_id}$.

\subsubsection{Transcript Authentication}
\label{sec:mint-dkg-transcript}

The full ceremony transcript $\mathcal{T}$ is the concatenation of all broadcast messages from Rounds 1--3, ordered canonically. The transcript hash is
\[
  \tau = \mathsf{SHA\text{-}256}(\mathsf{ceremony\_id} \,\|\, \mathsf{epoch} \,\|\, \mathcal{T}).
\]
Each validator $V_i$ in $\mathsf{QUAL}$ publishes a \emph{transcript signature}
\[
  \mathsf{TranscriptSignature}_i = \mathsf{Sign}(\mathsf{MIK}_i, \text{``tardus-ceremony-v1''} \,\|\, \tau)
\]
using its Tier 1 offline master key. The domain separator string distinguishes transcript signatures from all other signature usages in the protocol; the construction is a Schnorr signature over the prefixed message per~\Cref{con:schnorr}.

A ceremony is \emph{ratified} once at least $t$ valid transcript signatures from members of $\mathsf{QUAL}$ have been published on a canonical record (the on-chain ceremony registry, \Cref{sec:onchain}).

\subsection{Keyset Identifier Derivation}
\label{sec:mint-keysetid}

The mint operates a family of denomination keys; each denomination $d$ in the set of supported values $\mathcal{D} = \{2^0, 2^1, \ldots, 2^{32}\}$ lamports has its own joint key $PK_d$ produced by a separate DKG ceremony. The keyset identifier per~\Cref{sec:hash} and \texttt{research/PRODUCTION\_LESSONS.md} §L3 is
\[
  \mathsf{KeysetID}_d = \mathtt{0x02} \,\|\, \mathsf{SHA\text{-}256}\bigl( \mathsf{sorted\_pairs}(\mathcal{D}, \{PK_d\}_{d \in \mathcal{D}}) \,\|\, \text{``|unit:tardus-sol|''} \,\|\, \text{``|epoch:''} \,\|\, e \,\|\, \text{``|''} \bigr),
\]
where $e$ is the epoch number. The version byte $\mathtt{0x02}$ distinguishes TARDUS keyset IDs from Cashu v1 ($\mathtt{0x00}$) and Cashu v2 ($\mathtt{0x01}$). Clients are required to recompute this identifier locally upon receiving any mint response and reject mismatched identifiers (\Cref{sec:wallet}).

\subsection{Threshold Blind Signing}
\label{sec:mint-blindsig}

Given a ratified joint key $PK$ for a denomination $d$, the mint produces blind Schnorr signatures via the construction of~\Cref{con:blind-threshold}. The four-round protocol is realised concretely as follows. Throughout, $\mathsf{session\_id}$ is a fresh $128$-bit identifier sampled by the user, binding the four rounds together.

\subsubsection{Round 1 --- Commitment Aggregation}
\label{sec:mint-sign-r1}

The user submits an issuance request $(\mathsf{session\_id}, d)$ to the PAE. Each signing validator $V_i$ (in some chosen signing set $S$ with $|S| = t$, selected by stake-weighted rotation) samples $k_i \leftarrow_\$ \mathbb{F}_\ell$ inside its HSM, computes $R_i = k_i G$, and publishes $R_i$ to the aggregator (which may be the PAE of any participating validator). The aggregator computes
\[
  R = \sum_{i \in S} R_i
\]
and forwards $R$ to the user, along with the chosen signing set $S$ and a $\mathsf{TranscriptSignature}$ binding $(\mathsf{session\_id}, S, R)$ to the issuing committee.

\subsubsection{Round 2 --- User Blinding}
\label{sec:mint-sign-r2}

The user samples $\alpha, \beta \leftarrow_\$ \mathbb{F}_\ell$ and computes the unblinded commitment, blinded challenge, and forwards the latter:
\begin{align*}
  R' &= R + \alpha G + \beta \cdot PK \\
  c' &= H_{\mathbb{F}_\ell}(R' \,\|\, PK \,\|\, m) \\
  c &= c' + \beta \pmod{\ell}
\end{align*}
where $m$ is the coin payload (\Cref{sec:refresh}). The user retains $(\alpha, \beta, R', \mathsf{session\_id})$ as private state, persisted via the trust-boundary serialisation discipline of \Cref{sec:wallet}.

\subsubsection{Round 3 --- Partial Response Aggregation}
\label{sec:mint-sign-r3}

Each validator $V_i$ in $S$ computes its partial signature
\[
  s_i = k_i + c \cdot \lambda_i^{(S)} \cdot sk_i \pmod{\ell},
\]
where $\lambda_i^{(S)} = \prod_{j \in S, j \neq i} \frac{-j}{i - j} \pmod{\ell}$ is the Lagrange interpolation coefficient at $i$ over signing set $S$, and broadcasts $s_i$. The aggregator computes
\[
  s = \sum_{i \in S} s_i
\]
and forwards $s$ to the user.

\subsubsection{Round 4 --- User Unblinding}
\label{sec:mint-sign-r4}

The user computes
\[
  s' = s + \alpha \pmod{\ell}
\]
and emits the unblinded signature $\sigma' = (R', s')$, which by direct verification of $s' G = R' + c' \cdot PK$ is a valid Schnorr signature on $m$ under $PK$ per~\Cref{con:schnorr}.

\begin{remark}
\label{rem:nonce-reuse}
A validator $V_i$ that participates with the same $(k_i, \mathsf{session\_id})$ pair in two distinct rounds reveals its share $sk_i$ to anyone who observes both partial signatures (the standard Schnorr nonce-reuse attack, lifted to the threshold setting). The mint protocol prohibits this; each validator's HSM enforces a unique-$k$-per-$\mathsf{session\_id}$ invariant by tracking session IDs in a tamper-evident log. A validator caught violating this invariant is removed from the committee via the revoke path of \Cref{sec:mint-revoke}.
\end{remark}

\subsection{Proactive Secret-Sharing Rotation}
\label{sec:mint-rotation}

To bound long-term metadata-correlation risk and mitigate slow share-extraction attacks against individual validators, the joint secret is rotated every epoch ($1$ day) via a proactive secret-sharing protocol:

\begin{enumerate}[leftmargin=2em,itemsep=0.2em]
  \item At epoch boundary $e \to e + 1$, the committee initiates a \emph{reshare ceremony} with a fresh $\mathsf{ceremony\_id}$.
  \item Each validator $V_i$ acts as a VSS dealer for the zero polynomial: it samples $\tilde{f}_i(x) = \tilde{a}_{i,1} x + \cdots + \tilde{a}_{i,t-1} x^{t-1}$ with $\tilde{f}_i(0) = 0$. The shares $\tilde\sigma_{i \to j} = \tilde{f}_i(j)$ are distributed and verified analogously to~\Cref{sec:mint-dkg-r1}.
  \item Each validator $V_j$ updates its share:
    \[
      sk_j^{(e+1)} = sk_j^{(e)} + \sum_{i \in \mathsf{QUAL}^{(e+1)}} \tilde{f}_i(j) \pmod{\ell}.
    \]
  \item The joint public key $PK$ is unchanged because $\sum_i \tilde{f}_i(0) = 0$.
  \item Each validator securely erases $sk_j^{(e)}$ from HSM after $\mathsf{OVERLAP\_DURATION} = 2$ epochs to preserve in-flight signatures.
\end{enumerate}

\paragraph{Cheating-dealer protection.} Soundness against a reshare dealer that tries to shift the joint key requires the explicit verifier check $\tilde{A}_{i,0} = \mathcal{O}$ (the group identity), since $\tilde{A}_{i,0} = \tilde{f}_i(0) \cdot G$. A non-identity commitment indicates $\tilde{f}_i(0) \neq 0$, which would alter $PK$. Verifiers reject any reshare broadcast that fails this check.

After a reshare, an attacker who held $\geq t-1$ shares from epoch $e$ obtains no advantage in epoch $e+1$, because the share polynomial has changed despite the same $PK$.

The committee maintains a $\mathsf{LOOKAHEAD} = 30$-epoch buffer of pre-completed reshare ceremonies, so that scheduled or emergency rotation is always available without waiting for ceremony completion.

\subsection{Denomination Key Revocation}
\label{sec:mint-revoke}

If a denomination key is compromised (e.g., $\geq t$ validators' shares are extracted, or a critical implementation flaw is discovered), it must be revoked. Revocation is a threshold-signed on-chain transaction that:

\begin{enumerate}[leftmargin=2em,itemsep=0pt]
  \item Marks $PK_d$ as revoked in the on-chain keyset registry.
  \item Triggers the recoup path of~\Cref{sec:mint-recoup}.
  \item Schedules an emergency DKG ceremony for $d$ with a new $PK_d^{\text{new}}$.
\end{enumerate}

Because revocation requires $\geq t$ committee signatures, no individual operator can forcibly revoke a key. A revocation transaction structure is fully specified in~\Cref{sec:onchain}.

\subsection{Recoup Protocol}
\label{sec:mint-recoup}

Following a revocation, users holding coins issued under the revoked $PK_d$ exchange them for new coins under $PK_d^{\text{new}}$ via the recoup protocol:

\begin{enumerate}[leftmargin=2em,itemsep=0.2em]
  \item The user presents a coin $(s, \sigma)$ from the revoked keyset to the mint together with a fresh blind issuance request under $PK_d^{\text{new}}$.
  \item The mint verifies the old signature, marks the old serial as recouped (nullifying it), and proceeds with a normal blind issuance under $PK_d^{\text{new}}$.
  \item The user obtains a new coin of equal denomination.
\end{enumerate}

Recoup is rate-limited per epoch to prevent a denial-of-service campaign from saturating the mint; the rate is set such that all legitimate revoked coins can be migrated within $\mathsf{LOOKAHEAD}$ epochs.

\subsection{Slot-Windowed Batch Swap and Uniform Fee Enforcement}
\label{sec:mint-batch}

To prevent timing- and fee-based deanonymisation, TARDUS issuance transactions are restricted to specific slot windows during which all issuance transactions in flight are batched into a Jito bundle. Inside a bundle:

\begin{itemize}[leftmargin=2em,itemsep=0pt]
  \item Every TARDUS transaction pays a uniform protocol fee of $10^{-4}$ SOL.
  \item Priority-fee variation is prohibited; transactions with non-uniform priority fees are rejected by the on-chain program (\Cref{sec:onchain}).
  \item The bundle ordering is randomised by the bundle coordinator (chosen by stake-weighted rotation per epoch), masking submission order from the chain observer.
\end{itemize}

This design follows the mitigation analysis of round 4 (committee item R6) and addresses the MEV attack surface identified in \Cref{sec:security}.

\subsection{Fee-Payer Privacy Hardening}
\label{sec:mint-fee-payer-privacy}

The Solana base layer reveals the \texttt{fee\_payer} of every
transaction. If a TARDUS user always pays their own fees, an
on-chain graph crawler can correlate all of that user's spends to
a single wallet — regardless of how strong the off-chain blind sign
+ sealed-box stack is. v1.5 ships a three-layer mitigation:

\paragraph{Layer 1 — Ephemeral payer (Faz 9.1).} For each
spend-class TX (Refresh, Withdraw), the client generates a fresh
ed25519 keypair, funds it from a sponsor wallet with a small
amount (default $10^{-3}$ SOL), and uses the ephemeral as the TX
\texttt{fee\_payer}. The ephemeral keypair is then dropped;
remaining dust lamports stay stranded on-chain. Result: each
TARDUS spend has a previously-unseen signer.

\paragraph{Layer 2 — Multi-sponsor pool rotation (Faz 9.2).} The
single sponsor wallet of Layer 1 is replaced by a randomly-rotated
selection from a set of sponsor keypairs. Per-TX random sponsor
selection breaks the "all ephemeral payers funded from the same
sponsor" link.

\paragraph{Layer 3 — On-chain SponsorPool (Faz 9.3, 9.4, 9.5).}
A program-owned PDA pools SOL deposits from any number of
independent funders. \texttt{SponsorPayout} drains a fixed amount
(default $10^{-3}$ SOL) from the pool to a fresh ephemeral
keypair. Anyone can deposit; anyone can call payout; the pool
commingles. The funding chain becomes:
\[
  \{\text{anonymous depositors}\} \rightarrow \text{SponsorPool PDA} \rightarrow \text{fresh ephemeral} \rightarrow \text{Refresh TX}
\]
which decouples the original SOL source from the spend signer.
The pool is rate-limited (default: one payout per five Solana
slots, $\approx 2.5\,\text{s}$) to bound drain attacks.

\paragraph{Layer 4 — Pool caller rotation (Faz 9.5).} The wallet
that submits the \texttt{SponsorPayout} TX is itself selected at
random from the same multi-sponsor pool, so the
"deployer always calls payout" residual leak is closed.

\paragraph{Residual leaks not addressed by Layers 1--4.}

\begin{itemize}[leftmargin=2em,itemsep=0pt]
  \item The \texttt{SponsorPayout} TX is publicly visible (anyone
    can correlate "pool $\rightarrow$ ephemeral X" by reading the
    TX data). Mitigation: many ephemerals + many payouts $\rightarrow$
    statistical unlinkability via pool size. Production deployments
    should keep ephemeral-creation rate high enough that any
    particular ephemeral is one of many in flight.
  \item The TX submitter's IP address. Mitigation: Tor / I2P
    circuit between user wallet and Solana RPC endpoint
    (\Cref{sec:relay-failure-modes} R7, deployment concern).
  \item Per-slot timing fingerprinting. Mitigation: random
    submission delays + Jito bundle ordering
    (\Cref{sec:mint-batch}).
\end{itemize}

\paragraph{Threat model improvement (concrete numbers).} For 100
TARDUS spends over a week, with a sponsor pool of 5 wallets and
random pool caller selection:

\begin{center}
\small
\begin{tabular}{lrrr}
\toprule
Posture & Distinct signers & Distinct funders & Direct link rate \\
\midrule
Plain (v8): user pays own fee     & 1 (the user)         & 1                  & 100\% \\
Layer 1: ephemeral payer          & 100 (one per TX)     & 1 (single sponsor) & 100\% via funding chain \\
Layer 2: multi-sponsor pool       & 100                  & 5                  & $\approx 20\%$ per sponsor \\
Layer 3+4: on-chain SponsorPool   & 100                  & pool-aggregated    & 0\% direct, statistical only \\
\bottomrule
\end{tabular}
\end{center}

\subsection{Validator Daemon Reference Surface}
\label{sec:mint-validator-daemon}

The Phase 1 reference implementation of the mint-side committee
member is the \texttt{tardus-validator} daemon, a long-running HTTP
service that holds one share of the joint mint key, coordinates with
its $n-1$ peers over the ceremony protocols of
\Cref{sec:mint-dkg,sec:mint-blindsig,sec:mint-rotation}, and exposes
the user-facing signing endpoints. This subsection documents the
v2 endpoint surface; the wire-level Borsh encoding for each request
/ response payload lives in the daemon's
\texttt{crates/tardus-validator} source.

\paragraph{Ceremony coordination (peer-to-peer, mTLS).}

\begin{center}
\small
\begin{tabular}{ll}
\toprule
HTTP path & Role \\
\midrule
\verb|POST /dkg/start|       & each daemon's local share + broadcast \\
\verb|POST /dkg/contribute|  & cross-feed peer contributions \\
\verb|POST /dkg/finalize|    & finalise joint\_pk \\
\verb|POST /reshare/start|   & v2.6 zero-poly reshare initiation \\
\verb|POST /reshare/contribute| & peer reshare contribution intake \\
\verb|POST /reshare/finalize|   & rotate to new epoch's share \\
\bottomrule
\end{tabular}
\end{center}

\paragraph{User-facing signing (anonymous, TLS).}

\begin{center}
\small
\begin{tabular}{ll}
\toprule
HTTP path & Role \\
\midrule
\verb|POST /sign/round1|     & user's $R$ commitment intake; daemon returns $k$-nonce \\
\verb|POST /sign/round3|     & user's blinded $c$; daemon returns blind partial signature \\
\verb|POST /refresh/round1|  & user's $\{\hat{c}_\gamma\}$ commitments; daemon returns nonces \\
\verb|POST /refresh/round5|  & user's $\{c_\gamma\}$ after cut-and-choose reveal; daemon signs \\
\bottomrule
\end{tabular}
\end{center}

\paragraph{Operator probes (HTTP, no auth).}
\verb|GET /health| (uptime, ceremony counts, share count),
\verb|GET /info| (operator name, joint\_pk \emph{compressed}, bind
addr, share backend), \verb|GET /version| (build hash, dep
versions).

\paragraph{Share backend trait.} \texttt{tardus-validator} dispatches
share-record persistence through a \texttt{ShareStore} trait
implemented by:
\begin{itemize}[leftmargin=2em,itemsep=0pt]
  \item \texttt{FileShareStore} (v2.1) — AEAD-sealed
        \texttt{ChaCha20-Poly1305} ciphertext under
        \texttt{HKDF-SHA-512(master\_seed)}, atomic
        tmp-and-rename writes.
  \item \texttt{Pkcs11ShareStore} (v2.11, behind the \texttt{hsm}
        cargo feature) — HSM-resident AES-256-GCM wrap key with
        \texttt{CKA\_EXTRACTABLE=false}. The wrap key never leaves
        the device; the daemon round-trips through
        \texttt{C\_Encrypt}/\texttt{C\_Decrypt} per save/load.
\end{itemize}

\paragraph{Transparency log.} Every ceremony state transition
(DKG/sign/refresh/reshare per session) appends a hash-chained
JSON-lines event:
$\mathsf{event}_i = \{\mathsf{kind},\,\mathsf{ts},\,\mathsf{prev\_hash},\,\mathsf{payload\_digest}\}$
with $\mathsf{prev\_hash}_i = \mathsf{SHA\text{-}256}(\mathsf{event}_{i-1})$.
The genesis event has $\mathsf{prev\_hash} = \mathbf{0}^{32}$.
Operators publish the head hash hourly; an external auditor with the
head hash can reject any post-hoc rewrite (insertion, deletion,
reordering) of validator behaviour.

\subsection{Operational Failure Modes}
\label{sec:mint-failures}

A non-exhaustive enumeration of mint-side failure modes appears in~\Cref{sec:failure-modes}. The most operationally-relevant are:

\begin{itemize}[leftmargin=2em,itemsep=0pt]
  \item \textbf{Validator offline during signing.} If $|S \cap \mathsf{ONLINE}| < t$, the aggregator selects an alternative signing set from the lookahead schedule. If no $t$-subset is online, issuance is queued until the lookahead can be refreshed.
  \item \textbf{Ceremony abort.} A failed DKG or reshare ceremony is restarted with a fresh $\mathsf{ceremony\_id}$. Repeated failures (more than three within an epoch) trigger an alert and pause issuance under the affected denomination until manual review.
  \item \textbf{Nonce-reuse attempt.} As~\Cref{rem:nonce-reuse} describes, an HSM-enforced unique-$k$ invariant catches the canonical attack; a violating validator is automatically removed from the active committee via the revoke path of~\Cref{sec:mint-revoke}.
\end{itemize}

\section{Refresh Protocol}
\label{sec:refresh}

The refresh protocol is the primary spending and change-making operation of TARDUS. A user holding a valid coin under the current joint key surrenders it to the mint in exchange for one or more fresh coins whose serial numbers are unlinkable to the surrendered coin. Refresh subsumes both ``swap'' (one coin in, one coin out, same denomination) and ``change'' (one coin in, multiple coins out summing to the original denomination) in a single protocol.

The construction is the $\kappa$-fold cut-and-choose refresh of~\cite{dold2019} (DD-62), adapted to the threshold blind Schnorr setting of~\Cref{sec:mint-blindsig}. The default parameter is $\kappa = 3$, giving a per-attempt detection probability of $1 - 1/\kappa = 2/3$ against a cheating user.

\subsection{Coin Model}
\label{sec:refresh-coin}

A \emph{coin} is a triple $(x, Cp, \sigma)$:
\begin{itemize}[leftmargin=2em,itemsep=0pt]
  \item $x \in \mathbb{F}_\ell$ is the secret seed (the spending key).
  \item $Cp = x \cdot G$ is the coin's public commitment.
  \item $\sigma$ is a valid Schnorr signature of $Cp$ (in its 32-byte compressed encoding) under the joint mint key $PK$ for the coin's denomination.
\end{itemize}

To verify a coin, one checks $\mathsf{Verify}(PK, Cp, \sigma) = \mathbf{1}$ per~\Cref{con:schnorr}.

To spend a coin to the on-chain ledger (\Cref{sec:onchain}), the holder presents $(Cp, \sigma)$. The on-chain program verifies the signature against $PK$ and records the \emph{nullifier}
\[
  \mathsf{null}(Cp) = \mathsf{SHA\text{-}256}(\text{``TARDUS-nullifier-v1''} \,\|\, Cp.\mathsf{bytes}())
\]
in the compressed nullifier tree. Double-spend is prevented by uniqueness of nullifier insertion: any subsequent attempt to spend the same coin produces the same nullifier and is rejected.

\paragraph{Bearer-only authentication.} The coin secret $x$ is never revealed on-chain. Whoever holds $(Cp, \sigma)$ can spend the coin; the secret remains private (a wallet-only artefact retained for the user's bookkeeping). The earlier draft of this protocol tied the nullifier to $x$, which required the on-chain program to recompute $Cp = x \cdot G$ — incompatible with Solana SBF's stack budget (\texttt{research/PRODUCTION\_LESSONS.md} §R8, discovered in v1.4.2). The $\mathsf{null}(Cp)$ formulation moves all curve arithmetic off-chain to the Solana \texttt{ed25519} precompile.

\subsection{Session Parameters}
\label{sec:refresh-params}

A refresh session is parameterised by:
\begin{itemize}[leftmargin=2em,itemsep=0pt]
  \item $\mathsf{session\_id} \in \{0,1\}^{128}$ --- fresh per session, prevents cross-session replay and binds the four threshold-signing rounds (\Cref{sec:mint-blindsig}) together within the refresh.
  \item $\kappa \in \mathbb{N}$ --- cut-and-choose parameter; default $\kappa = 3$.
  \item $S \subseteq [n]$ --- signing set of size $t$ chosen from the active committee via stake-weighted rotation.
  \item $\mathsf{old\_coin} = (x_{\mathsf{old}}, Cp_{\mathsf{old}}, \sigma_{\mathsf{old}})$ --- the coin being surrendered.
\end{itemize}

\subsection{Coin Secret Derivation}
\label{sec:refresh-derivation}

Fresh coin secrets are derived from per-candidate seeds via HKDF, never from a long-term wallet seed via BIP32. This is the explicit lesson from the Cashu NUT-13 disclosure (\texttt{research/PRODUCTION\_LESSONS.md} §L2): deterministic seed-based scalar derivation creates the keyset-collision attack surface; we forbid it.

Concretely, the user samples a fresh $\mathsf{seed}_\gamma \in \{0,1\}^{256}$ uniformly for each candidate $\gamma \in [1, \kappa]$, and derives:
\begin{align*}
  \mathsf{ikm}_\gamma   &\gets \mathsf{seed}_\gamma, \\
  \mathsf{prk}_\gamma   &\gets \mathsf{HKDF\text{-}Extract}(\text{``TARDUS-refresh-v1''}, \mathsf{ikm}_\gamma), \\
  \mathsf{okm}_\gamma   &\gets \mathsf{HKDF\text{-}Expand}(\mathsf{prk}_\gamma, \text{``coin-secret''}, 64\text{ bytes}), \\
  x_\gamma            &\gets \mathsf{okm}_\gamma \bmod \ell,
\end{align*}
with the reduction-bias bound $\ell / 2^{512} \approx 2^{-260}$ (negligible; cf.~\Cref{sec:hash}). The seed is one-shot and never reused, so the deterministic-collision surface that NUT-13 exposed cannot arise.

\subsection{Six-Round Protocol}
\label{sec:refresh-rounds}

The refresh protocol runs six rounds. Rounds 1, 3, and 5 are mint-initiated; Rounds 2, 4, and 6 are user-initiated. Each round message carries the $\mathsf{session\_id}$ in the clear; mismatches abort the session (\Cref{sec:refresh-abort}).

\subsubsection{Round 1 --- Mint \texorpdfstring{$\kappa$}{kappa}-fold Commitment}
\label{sec:refresh-r1}

Each signing validator $V_i \in S$ samples $\kappa$ independent nonces $k_{i,\gamma} \leftarrow_\$ \mathbb{F}_\ell$ for $\gamma \in [1, \kappa]$, computes $R_{i,\gamma} = k_{i,\gamma} \cdot G$, and broadcasts the set. The aggregator computes
\[
  R_\gamma = \sum_{i \in S} R_{i,\gamma} \quad \text{for each } \gamma \in [1, \kappa]
\]
and forwards the $\kappa$ aggregated points to the user.

\paragraph{Nonce-reuse invariant.} The HSM-enforced unique-$(k_{i,\gamma}, \mathsf{session\_id}, \gamma)$ tracking of~\Cref{sec:mint-blindsig} extends to the $\kappa$-fold case: $\gamma$ is included in the tracking tuple.

\subsubsection{Round 2 --- User Blinded Challenges}
\label{sec:refresh-r2}

The user samples $\kappa$ fresh seeds, derives candidate coin secrets per~\Cref{sec:refresh-derivation}, and produces $\kappa$ blinded challenges. For each $\gamma \in [1, \kappa]$:
\begin{enumerate}[leftmargin=2em,itemsep=0.1em]
  \item Sample $\mathsf{seed}_\gamma \leftarrow_\$ \{0,1\}^{256}$, derive $x_\gamma$.
  \item Compute $Cp_\gamma = x_\gamma \cdot G$, encode $\mathsf{msg}_\gamma = Cp_\gamma.\mathsf{compress}().\mathsf{bytes}()$.
  \item Sample $\alpha_\gamma, \beta_\gamma \leftarrow_\$ \mathbb{F}_\ell$.
  \item Compute the blinded commitment $R'_\gamma = R_\gamma + \alpha_\gamma G + \beta_\gamma \cdot PK$.
  \item Compute the blinded challenge $c_\gamma = H_{\mathbb{F}_\ell}(R'_\gamma \,\|\, PK \,\|\, \mathsf{msg}_\gamma) + \beta_\gamma \pmod{\ell}$.
\end{enumerate}
The user submits $\{c_\gamma\}_{\gamma=1}^{\kappa}$ to the mint together with the surrendered $\mathsf{old\_coin}$.

\subsubsection{Round 3 --- Mint Cut-and-Choose}
\label{sec:refresh-r3}

The mint verifies $\mathsf{Verify}(PK, Cp_{\mathsf{old}}, \sigma_{\mathsf{old}}) = \mathbf{1}$. If valid, the mint samples $\gamma^\star \leftarrow_\$ [1, \kappa]$ uniformly at random (via the committee VRF; deterministic from $\mathsf{session\_id}$ to avoid manipulation) and sends $\gamma^\star$ to the user as the challenge.

\subsubsection{Round 4 --- User Reveal}
\label{sec:refresh-r4}

The user reveals $\{ (\mathsf{seed}_\gamma, \alpha_\gamma, \beta_\gamma) \}_{\gamma \neq \gamma^\star}$ --- all $\kappa - 1$ candidates except the challenged one. The mint, for each revealed $\gamma$, re-derives $x_\gamma$, $Cp_\gamma$, $R'_\gamma$, recomputes $c_\gamma$, and verifies that the recomputed value matches the $c_\gamma$ the user submitted in Round 2.

If any verification fails, the session aborts (\Cref{sec:refresh-abort}) and the old coin's nullifier is~\emph{not} inserted; the user must retry from Round 1 with a fresh session.

\subsubsection{Round 5 --- Mint Partial Signing}
\label{sec:refresh-r5}

If all $\kappa-1$ revealed candidates verified, the mint proceeds to sign the $\gamma^\star$-th candidate. Each $V_i \in S$ computes the partial signature
\[
  s_{i, \gamma^\star} = k_{i, \gamma^\star} + c_{\gamma^\star} \cdot \lambda_i^{(S)} \cdot sk_i \pmod{\ell}
\]
exactly as in~\Cref{sec:mint-sign-r3}. The aggregator returns
\[
  s_{\gamma^\star} = \sum_{i \in S} s_{i, \gamma^\star}
\]
to the user, together with proof that $Cp_{\mathsf{old}}$ has been inserted into the nullifier tree (the on-chain transaction signature).

\subsubsection{Round 6 --- User Unblinding}
\label{sec:refresh-r6}

The user computes $s'_{\gamma^\star} = s_{\gamma^\star} + \alpha_{\gamma^\star} \pmod{\ell}$. The new coin is
\[
  (x_{\gamma^\star}, Cp_{\gamma^\star}, \sigma_{\gamma^\star}) \quad \text{where} \quad \sigma_{\gamma^\star} = (R'_{\gamma^\star}, s'_{\gamma^\star}).
\]
By construction, $\mathsf{Verify}(PK, Cp_{\gamma^\star}, \sigma_{\gamma^\star}) = \mathbf{1}$, and the new coin's $Cp_{\gamma^\star}$ is statistically unlinkable to $Cp_{\mathsf{old}}$ (no validator has seen $x_{\gamma^\star}$ or the unblinding factors).

\subsection{Anonymity Argument}
\label{sec:refresh-anon}

For a malicious mint coalition of size $< t$, the only data observed during refresh are:
\begin{enumerate}[leftmargin=2em,itemsep=0pt]
  \item The surrendered $(Cp_{\mathsf{old}}, \sigma_{\mathsf{old}})$ and the user's challenges $\{c_\gamma\}$.
  \item The challenge $\gamma^\star$ (publicly fixed by VRF).
  \item The revealed candidates' $(seed_\gamma, \alpha_\gamma, \beta_\gamma)$ for $\gamma \neq \gamma^\star$.
  \item The partial signature contributions $\{s_{i, \gamma^\star}\}_{i \in S}$.
\end{enumerate}
The new coin's secret $x_{\gamma^\star}$ and pubkey $Cp_{\gamma^\star}$ are not in this view. The blinded challenge $c_{\gamma^\star}$ is, by the statistical-blindness argument of~\Cref{thm:blindness}, statistically independent of the unblinded signature $\sigma_{\gamma^\star}$ and message $Cp_{\gamma^\star}$. Therefore, the mint coalition cannot link the new coin to the old coin with probability exceeding $1/2$ over the natural sample space.

\subsection{Cheating-User Detection Rate}
\label{sec:refresh-detection}

A cheating user who substitutes a $c_\gamma$ that is inconsistent with their announced derivation must hope that the mint chooses $\gamma^\star = \gamma$ (the inconsistent candidate). Otherwise, the inconsistency is exposed by the Round-4 verification. Across an i.i.d.~uniform random $\gamma^\star$, the probability of evasion is $1/\kappa$. For $\kappa = 3$, this is $33\%$ per attempt --- a single attempt is detected with $67\%$ probability, two attempts with $89\%$, three with $96\%$. Repeated cheating attempts thus give the mint information-theoretic confidence in malicious user identification, while leaving an honest user with zero detection probability.

\subsection{Abort Handling}
\label{sec:refresh-abort}

Any of the following triggers a session abort:
\begin{itemize}[leftmargin=2em,itemsep=0pt]
  \item $\mathsf{session\_id}$ mismatch in any round message.
  \item Invalid $\sigma_{\mathsf{old}}$ in Round 3.
  \item Failed reveal verification in Round 4.
  \item Threshold signing failure in Round 5 (e.g., $< t$ validators online).
\end{itemize}

On abort:
\begin{enumerate}[leftmargin=2em,itemsep=0pt]
  \item The old coin's nullifier is~\emph{not} inserted; the coin remains spendable.
  \item All session state (per-validator $k_{i,\gamma}$, user $\alpha_\gamma, \beta_\gamma, \mathsf{seed}_\gamma$) is securely erased.
  \item The user may retry with a fresh $\mathsf{session\_id}$.
\end{enumerate}

The user's recoverable state across a Round-1-through-Round-6 protocol fault is captured by serialising the user's $\kappa$-tuple $\{\mathsf{seed}_\gamma, \alpha_\gamma, \beta_\gamma, R'_\gamma, c_\gamma\}$ to encrypted local storage between Rounds 2 and 6 (\Cref{sec:wallet}).

\section{On-Chain Solana Program}
\label{sec:onchain}

The on-chain component of TARDUS is a single native Solana program that custodies the collateral vault, maintains the keyset registry, enforces the nullifier set, and verifies coin spends. The program is written in pure Rust (no Anchor) and runs on the Solana BPF / SBF virtual machine. This section specifies the account model, the instruction set, the per-instruction compute-unit budget, and the security invariants the program enforces.

\subsection{Program Overview}
\label{sec:onchain-overview}

The on-chain program is the trust anchor between the off-chain threshold mint (\Cref{sec:mint}) and the public Solana ledger. It performs no signing of its own --- all coin signatures originate off-chain from the threshold committee. Its responsibilities are:

\begin{itemize}[leftmargin=2em,itemsep=0pt]
  \item Verify coin signatures produced off-chain (via the Solana \texttt{ed25519} signature verification path).
  \item Maintain the nullifier tree, preventing double-spend (\Cref{sec:refresh-coin}).
  \item Hold the SPL Token-2022 vault that backs all issued coins of a denomination.
  \item Track the active keyset registry so verifiers know which $PK$ to use per denomination per epoch.
  \item Enforce the slot-windowed batch and uniform-fee discipline of~\Cref{sec:mint-batch} when accepting refresh transactions.
\end{itemize}

The program is \emph{passive}: it does not initiate any operation. All instructions are submitted by users or by the threshold committee (for revoke/registry maintenance).

\subsection{Account Model}
\label{sec:onchain-accounts}

\subsubsection{Keyset Registry}
\label{sec:onchain-keyset-registry}

A single program-derived account holds the canonical list of active keysets:
\[
  \mathsf{KeysetRegistry} = (\mathsf{version}, [\mathsf{KeysetEntry}_1, \ldots, \mathsf{KeysetEntry}_K])
\]
with
\[
  \mathsf{KeysetEntry} = (\mathsf{keyset\_id} \in \{0,1\}^{264}, \mathsf{denom} \in \{2^0, \ldots, 2^{32}\}, \mathsf{joint\_pk} \in \{0,1\}^{256}, \mathsf{epoch} \in \mathbb{N}, \mathsf{status} \in \{\mathsf{Active}, \mathsf{Revoked}\}).
\]
The PDA seed is $(\text{``tardus''}, \text{``keyset-registry''}, \text{program\_id})$. The account is initialised once at program deployment by a designated bootstrap authority and subsequently updated only via threshold-signed instructions (\Cref{sec:onchain-instr-register}).

\subsubsection{Vault}
\label{sec:onchain-vault}

For each denomination $d \in \mathcal{D}$, the program owns a Token-2022 Confidential Mint vault account:
\[
  \mathsf{Vault}_d = \mathsf{Token\text{-}2022\ Confidential\ Mint}(\mathsf{authority}=\mathsf{PDA}(d), \mathsf{denom}=d)
\]
with the PDA seed $(\text{``tardus''}, \text{``vault''}, d.\mathsf{to\_le\_bytes}(), \text{program\_id})$.

The collateral invariant (spec §3.10, EasyCrypt T4) is:
\[
  \forall \text{ slot } s: \mathsf{Vault}_d.\mathsf{collateral}(s) = \sum_{cm \in \mathsf{active\_commitments}_d(s)} d.
\]
This invariant is preserved atomically across every \texttt{Deposit}, \texttt{Withdraw}, \texttt{Refresh}, and \texttt{Recoup} instruction. The program rejects any state transition that would violate it.

\subsubsection{Nullifier Tree}
\label{sec:onchain-nullifier-tree}

Spent coins are recorded by their nullifier $\mathsf{null}(x) = \mathsf{SHA\text{-}256}(\text{``TARDUS-nullifier-v1''} \,\|\, x)$ in a Light Protocol compressed Merkle tree of depth $d = 32$ (capacity $\approx 4.3$~billion leaves). The PDA seed is $(\text{``tardus''}, \text{``nullifier-tree''}, \text{program\_id})$.

Each \texttt{Refresh} or \texttt{Withdraw} instruction inserts the surrendered coin's nullifier; duplicate insertion (double-spend) is rejected by the Light account-compression CPI.

\subsection{Instruction Set}
\label{sec:onchain-instructions}

\subsubsection{Bootstrap}
\label{sec:onchain-instr-bootstrap}

Allocates a program-owned PDA at canonical seeds. The deployer pays
rent exemption from their own lamports; the program uses
\verb|invoke_signed(SystemInstruction::CreateAccount, ...)| under the
canonical PDA seeds + bump to create the account on its own behalf.
Idempotent: a second Bootstrap against an existing PDA returns
\verb|ProgramError::Custom(15)| (\texttt{ERR\_ACCOUNT\_ALREADY\_EXISTS}).

\paragraph{Inputs.}
\begin{itemize}[leftmargin=2em,itemsep=0pt]
  \item $\mathsf{account\_kind} \in \{\mathsf{KeysetRegistry}, \mathsf{NullifierTree}, \mathsf{Vault}\}$.
  \item $\mathsf{size}$ (u32) --- allocation size in bytes
        (ignored for $\mathsf{Vault}$, which is system-owned and
        lamports-only). Capped at $64\,\mathrm{KiB}$ to prevent DoS.
  \item $\mathsf{denom}$ (u64) --- denomination for $\mathsf{Vault}$ kind;
        ignored otherwise.
\end{itemize}

\paragraph{Effect.} Creates a new program-owned (or system-owned for
$\mathsf{Vault}$) PDA at the canonical seeds for the kind. Subsequent
\texttt{RegisterKeyset}, \texttt{Refresh}, \texttt{Withdraw}
instructions require these PDAs to already exist.

\paragraph{Discovery.} Added in v1.4.7 after early devnet integration
showed that pre-population of program-owned PDAs via off-chain tools
is not possible: only the program itself can own the data accounts it
needs, and \verb|invoke_signed| is the only path to materialise them
on-chain. \texttt{tardus-cli devnet bootstrap-singletons} and
\texttt{tardus-cli devnet bootstrap-vault --denom N} are the
operator-facing front of this instruction.

\subsubsection{RegisterKeyset}
\label{sec:onchain-instr-register}

Inserts a new keyset entry into the registry. Permitted only via a transaction signed by at least $t$ of the threshold committee's offline master keys.

\paragraph{Inputs.}
\begin{itemize}[leftmargin=2em,itemsep=0pt]
  \item $\mathsf{keyset\_id}$ ($33$ bytes, derived per~\Cref{sec:mint-keysetid}).
  \item $\mathsf{denom}$ (u64).
  \item $\mathsf{joint\_pk}$ ($32$ bytes).
  \item $\mathsf{epoch}$ (u64).
  \item $[\mathsf{TranscriptSignature}_1, \ldots, \mathsf{TranscriptSignature}_t]$ (ceremony ratification per~\Cref{sec:mint-dkg-transcript}).
\end{itemize}

\paragraph{Checks.}
\begin{itemize}[leftmargin=2em,itemsep=0pt]
  \item Each transcript signature verifies against a registered validator MIK.
  \item Submitted $\mathsf{keyset\_id}$ matches the canonical derivation from the new $\mathsf{joint\_pk}$.
  \item Registry has space for the new entry.
\end{itemize}

\paragraph{Effect.} Appends a new \texttt{Active} entry to the registry. Emits a Solana log event.

\subsubsection{Deposit}
\label{sec:onchain-instr-deposit}

User sends SOL (or SPL token) to the program's vault for a given denomination, in exchange for an off-chain mint-issued coin.

\paragraph{Inputs.}
\begin{itemize}[leftmargin=2em,itemsep=0pt]
  \item $\mathsf{denom}$ (u64).
  \item Lamports transferred via system transfer instruction in the same atomic transaction.
\end{itemize}

\paragraph{Checks.}
\begin{itemize}[leftmargin=2em,itemsep=0pt]
  \item Denomination is registered and \texttt{Active}.
  \item Transferred amount equals $\mathsf{denom}$ exactly.
\end{itemize}

\paragraph{Effect.} Credits the vault for the denomination. Off-chain, the user proceeds with the threshold mint protocol (\Cref{sec:mint-blindsig}) to receive a coin of the same denomination. The deposit instruction logs the deposit event; the actual coin issuance is mediated by the off-chain mint committee.

\subsubsection{Refresh}
\label{sec:onchain-instr-refresh}

Atomically: verify the surrendered coin's signature, insert its nullifier, and record the issuance of one or more new coins of equal total denomination.

\paragraph{Inputs.}
\begin{itemize}[leftmargin=2em,itemsep=0pt]
  \item $Cp_{\mathsf{old}}$ ($32$ bytes) --- the surrendered coin's public commitment.
  \item $\mathsf{denom}$ (u64) --- denomination of the surrendered coin.
  \item Off-chain prerequisite: a preceding \texttt{ed25519\_program} precompile instruction in the same transaction, verifying $\mathsf{Verify}(\mathsf{joint\_pk}, Cp_{\mathsf{old}}, \sigma_{\mathsf{old}}) = \mathbf{1}$. The \texttt{Refresh} instruction reads the instructions sysvar to confirm the precompile was invoked with the expected $(\mathsf{joint\_pk}, Cp_{\mathsf{old}})$ pair.
  \item The new coins $(Cp_{\mathsf{new}}, \sigma_{\mathsf{new}})$ emerge off-chain via the round-6 unblinding of~\Cref{sec:refresh}; they are \emph{not} on-chain in v1.4.3 (anonymity preservation).
\end{itemize}

\paragraph{Checks.}
\begin{itemize}[leftmargin=2em,itemsep=0pt]
  \item $\mathsf{denom}$ keyset is \texttt{Active}.
  \item Precompile pre-instruction verified $\sigma_{\mathsf{old}}$ against $(\mathsf{joint\_pk}, Cp_{\mathsf{old}})$.
  \item $\mathsf{null}(Cp_{\mathsf{old}}) = \mathsf{SHA\text{-}256}(\text{``TARDUS-nullifier-v1''} \,\|\, Cp_{\mathsf{old}})$ is not already in the nullifier tree.
\end{itemize}

\paragraph{Effect.} Inserts $\mathsf{nullifier}$ into the nullifier tree. Logs each new coin's issuance (without revealing the coin's secret).

\subsubsection{Withdraw}
\label{sec:onchain-instr-withdraw}

User redeems a coin for the underlying SOL (or SPL token).

\paragraph{Inputs.} Mirrors the v1.4.3 bearer model used by
\texttt{Refresh}:
\begin{itemize}[leftmargin=2em,itemsep=0pt]
  \item $Cp$ ($32$ bytes) --- the surrendered coin's public commitment.
  \item $\mathsf{denom}$ (u64).
  \item $\mathsf{recipient}$ ($32$ bytes Pubkey) --- destination for
        the released lamports.
  \item Off-chain prerequisite: same ed25519 precompile pre-instruction
        pattern as \texttt{Refresh} (\Cref{sec:onchain-instr-refresh}).
\end{itemize}

\paragraph{Checks.}
\begin{itemize}[leftmargin=2em,itemsep=0pt]
  \item $\mathsf{denom}$ keyset is \texttt{Active}.
  \item Precompile pre-instruction verified $\sigma$ against $(\mathsf{joint\_pk}, Cp)$.
  \item $\mathsf{null}(Cp)$ not already present in the nullifier tree.
  \item Vault PDA has $\geq \mathsf{denom}$ lamports.
  \item Recipient account matches the instruction's $\mathsf{recipient}$ field.
\end{itemize}

\paragraph{Effect.} Inserts nullifier; releases $\mathsf{denom}$
lamports from the vault PDA to the recipient via
\verb|invoke_signed(SystemInstruction::Transfer, vault → recipient)|
under the vault's canonical PDA seeds.

\subsubsection{Revoke}
\label{sec:onchain-instr-revoke}

Permits the threshold committee to mark a keyset as revoked. Required for compromise response (\Cref{sec:mint-revoke}).

\paragraph{Inputs.}
\begin{itemize}[leftmargin=2em,itemsep=0pt]
  \item $\mathsf{keyset\_id}$ to revoke.
  \item $[\mathsf{TranscriptSignature}_1, \ldots, \mathsf{TranscriptSignature}_t]$ (threshold-signed authorization).
\end{itemize}

\paragraph{Effect.} Updates the keyset entry's $\mathsf{status}$ to \texttt{Revoked}. Subsequent \texttt{Refresh} and \texttt{Withdraw} instructions against this keyset use the \texttt{Recoup} path (deferred to v1.5).

\subsubsection{SponsorDeposit / SponsorPayout (v1.4.13)}
\label{sec:onchain-instr-sponsor}

The on-chain \texttt{SponsorPool} PDA is a community-funded SOL
faucet that decouples spend-class TX signers from any specific
funding source (cross-reference \Cref{sec:mint-fee-payer-privacy}).

\paragraph{SponsorDeposit.}
\begin{description}[leftmargin=2em,itemsep=0pt]
  \item[Inputs] $\mathsf{amount}$ (u64).
  \item[Accounts] funder (signer), \texttt{SponsorPool} PDA, system program.
  \item[Effect] The TX MUST include a paired \texttt{system\_instruction::transfer}
    from \texttt{funder} to the \texttt{SponsorPool} PDA in the same TX (atomic).
    The handler updates the running \texttt{total\_deposits} counter.
\end{description}

\paragraph{SponsorPayout.}
\begin{description}[leftmargin=2em,itemsep=0pt]
  \item[Inputs] $\mathsf{lamports}$ (u64), $\mathsf{recipient}$ (Pubkey).
  \item[Accounts] caller (signer), \texttt{SponsorPool} PDA, recipient, system program.
  \item[Checks]
    \begin{itemize}[leftmargin=2em,itemsep=0pt]
      \item Rate limit: $\mathsf{current\_slot} - \mathsf{last\_payout\_slot} \geq 5$ (otherwise \texttt{ERR\_RATE\_LIMIT = 30}).
      \item Pool has $\geq \mathsf{lamports}$ available.
    \end{itemize}
  \item[Effect] Direct lamport movement (program-owned PDA debit) from
    \texttt{SponsorPool} to \texttt{recipient}. Updates \texttt{last\_payout\_slot}
    and \texttt{total\_payouts}.
\end{description}

\paragraph{Drain bound.} A pure-drain attacker (no honest deposits)
can withdraw at most $\mathsf{lamports} / (5 \cdot 400\,\text{ms}) =
0.5 \cdot \mathsf{lamports}\,\text{per second}$. With
$\mathsf{lamports} = 10^6$ (default 0.001 SOL/payout), the steady-state
drain ceiling is $\approx 14\,\text{SOL/hour}$. Honest deposit cadence
must exceed this for the pool to remain liquid.

\subsubsection{ResizeAccount (v1.4.14, scaling)}
\label{sec:onchain-instr-resize}

Resizes any program-owned PDA (\texttt{KeysetRegistry}, \texttt{NullifierTree},
\texttt{SponsorPool}) to a larger byte allocation. Required because
the initial \texttt{Bootstrap} allocations cap at $\approx 11$
keysets (1024 B) or $\approx 250$ nullifiers (8 KiB).

\paragraph{Inputs.}
\begin{itemize}[leftmargin=2em,itemsep=0pt]
  \item $\mathsf{account\_kind} \in \{\mathsf{KeysetRegistry}, \mathsf{NullifierTree}, \mathsf{SponsorPool}\}$.
  \item $\mathsf{new\_size}$ (u32, cap 64 KiB).
\end{itemize}

\paragraph{Checks.}
\begin{itemize}[leftmargin=2em,itemsep=0pt]
  \item PDA matches canonical seeds.
  \item $\mathsf{new\_size} >$ current size.
  \item Account lamports $\geq \mathsf{rent.minimum\_balance(new\_size)}$;
    the TX MUST include a paired \texttt{system\_instruction::transfer}
    for the rent top-up.
\end{itemize}

\paragraph{Effect.} \texttt{AccountInfo::realloc(new\_size, true)} —
zero-initialises the new tail bytes so subsequent Borsh
deserialisation with trailing-padding tolerance
(\Cref{sec:onchain-security} strict-deserialisation invariant) still
parses correctly.

\paragraph{Solana realloc 10 KB per-TX limit (discovery \#10).}
The Solana SBF runtime enforces \texttt{MAX\_PERMITTED\_DATA\_INCREASE
= 10240 bytes per instruction}. Growing from 1024 B to 65536 B
therefore requires $\lceil (65536 - 1024)/10240 \rceil = 7$ sequential
\texttt{ResizeAccount} calls. The CLI's
\texttt{tardus devnet resize-nullifier-tree} subcommand handles the
multi-TX batching transparently; operators see a single high-level
target size.

Long-term replacement: Light Protocol compressed-Merkle-tree
(\texttt{deploy/runbooks/light-protocol-integration-design.md}, v1.5+).

\subsection{Compute-Unit Budget}
\label{sec:onchain-cu}

\paragraph{Measured BPF runtime cost (v1.4.6).} The Phase 0 estimates
shipped in the original draft of this section were 3--10$\times$
higher than the runtime measurements obtained after the SBF binary
was deployed on Solana devnet. The current numbers, captured by
\texttt{tests/sbf\_integration.rs} via \texttt{solana-program-test}
running in BPF mode (the same SVM Solana uses):

\begin{center}
\begin{tabular}{lrr}
\toprule
Instruction & Measured CU & Headroom vs Phase 0 estimate \\
\midrule
Bootstrap (each kind)             & $7{,}500 - 9{,}000$ & --- (new in v1.4.7) \\
RegisterKeyset                    & $13{,}441$          & $\sim 7\times$ better \\
Refresh (1 in, ed25519 precompile bridge) & $67{,}558$  & $\sim 4\times$ better \\
Withdraw (precompile + invoke\_signed CPI) & $77{,}152$ & $\sim 3\times$ better \\
Refresh-reject (precompile missing) & $4{,}840$ before abort & --- \\
\bottomrule
\end{tabular}
\end{center}

Per Solana's $1{,}400{,}000$ CU per-transaction limit, the most
expensive single instruction (\texttt{Withdraw}) uses approximately
$5.5\%$ of the budget. The default $200{,}000$ CU budget per
instruction --- without an explicit \texttt{ComputeBudget} prelude ---
also accommodates every TARDUS instruction comfortably; no caller
needs to bump the limit.

\paragraph{Architectural shift documented.} The Phase 0 estimate
assumed in-program \texttt{ed25519} signature verification via
\texttt{curve25519-dalek} ($\sim 5{,}000$ CU/sig). Implementation
discovered (\texttt{research/PRODUCTION\_LESSONS.md} §R8) that
\texttt{curve25519-dalek}'s variable-base scalar-mul lookup-table
allocation exceeds the 4 KiB SBF stack budget. The fix
(\Cref{sec:onchain-instr-refresh} and v1.4.4) delegates signature
verification to Solana's \texttt{ed25519\_program} precompile via the
instructions-sysvar introspection pattern; the in-program cost is now
the bridge check ($\sim 6{,}000$ CU) rather than the verification
itself. The precompile's own cost ($\sim 3{,}000$ CU) is paid by the
caller in a preceding instruction in the same transaction.

\paragraph{Light Protocol deferred.} The Phase 0 estimate included
$\sim 40{,}000$ CU for Light Protocol validity-proof verification.
v1 ships with a plain Borsh-encoded \texttt{NullifierSet} (capacity
$\sim 250$ entries before the 8 KiB account fills); the Light
Protocol compressed-Merkle-tree integration is deferred to v1.4.8
(\Cref{sec:failure-modes} F14).

\subsection{PDA Derivation}
\label{sec:onchain-pda}

All program-owned accounts use Program-Derived Addresses with the following canonical seed schemes:

\begin{center}
\begin{tabular}{lll}
\toprule
Account & PDA seed bytes & Owner \\
\midrule
KeysetRegistry & \texttt{("tardus", "keyset-registry")}      & program           \\
NullifierTree  & \texttt{("tardus", "nullifier-tree")}       & program           \\
Vault for $d$  & \texttt{("tardus", "vault", d.le\_bytes())} & system program    \\
\bottomrule
\end{tabular}
\end{center}

\paragraph{Vault ownership rationale.} Vault PDAs are
\emph{system-owned} (not program-owned), even though the program
controls them. This is required by Solana semantics: only the system
program can debit lamports from accounts it owns, so the vault must
be system-owned for \texttt{Withdraw}'s
\verb|invoke_signed(SystemInstruction::Transfer)| to debit it.
Program-owned data PDAs (\texttt{KeysetRegistry},
\texttt{NullifierTree}) hold no lamports beyond rent exemption and
are debited only via account closure (out of scope for v1).

The program enforces canonical-PDA validation on every instruction
(\texttt{research/PRODUCTION\_LESSONS.md} §L4): each account argument is
checked to be the canonical PDA derivable from these seeds. This
blocks the account-confusion attack class
(\texttt{research/PRODUCTION\_LESSONS.md} §R7).

\subsection{Devnet Deployment Anchor}
\label{sec:onchain-deployment}

The TARDUS reference program is deployed on Solana devnet for
ongoing integration validation:

\begin{center}
\small
\begin{tabular}{ll}
\toprule
Field & Value \\
\midrule
Program ID            & \texttt{AmY1ysgQyCC6CmorXkrNkogBSHmomy4kbsPziAYUx47u} \\
Authority             & \texttt{5EiLNkpp3QiH1wzBEHnCdWSpxREFH2q5jje9S2gVmaMz} \\
Loader                & \texttt{BPFLoaderUpgradeab1e11111111111111111111111}    \\
Network               & devnet (\texttt{https://api.devnet.solana.com}) \\
Binary size           & $200\,\mathrm{KiB}$ ELF (\texttt{target/deploy/tardus\_program.so}) \\
Rent-locked           & $\approx 1.40\,\mathrm{SOL}$                                \\
\bottomrule
\end{tabular}
\end{center}

Any party with a Solana wallet may submit \texttt{Bootstrap},
\texttt{RegisterKeyset}, \texttt{Deposit}, \texttt{Refresh},
\texttt{Withdraw}, or \texttt{Revoke} transactions against this
program ID; the protocol is exercised continuously by
\texttt{tests/devnet\_e2e.rs} (marked \texttt{\#[ignore]} for default
\texttt{cargo test} runs, gated behind explicit \texttt{cargo test
-{}- -{}-ignored}).

\subsection{Security Invariants}
\label{sec:onchain-security}

The program enforces the following invariants at the boundary of every instruction:

\begin{itemize}[leftmargin=2em,itemsep=0pt]
  \item \textbf{Canonical PDA}: each account argument is checked to be the canonical PDA derivable from \Cref{sec:onchain-pda}'s seeds and the program ID.
  \item \textbf{Signer verification}: instructions that authenticate to the committee (\texttt{RegisterKeyset}, \texttt{Revoke}) check that at least $t$ valid committee transcript signatures are present.
  \item \textbf{Vault collateral invariant}: any instruction that changes vault collateral or the active-commitment set checks the invariant pre- and post-execution.
  \item \textbf{Strict deserialization with padding tolerance}: instruction data is parsed via Borsh \texttt{try\_from\_slice} (which rejects trailing bytes); account-data deserialization uses \texttt{BorshDeserialize::deserialize\_reader} (which reads only the encoded prefix and tolerates trailing zero padding). The reader-based path was added in v1.4.5 after a real-runtime SVM bug caused account-data loads to fail when the program-owned PDA was larger than the serialised state.
  \item \textbf{No CPI reentrancy}: the program does not recursively invoke itself; any CPI is to allow-listed external programs (Token-2022, Light account-compression).
  \item \textbf{Sysvar conservatism}: epoch transitions are derived from the Clock sysvar but tolerated within a slot-window buffer to absorb leader slot skew.
\end{itemize}

\section{Wallet Layer}
\label{sec:wallet}

The wallet layer is the user-facing surface of TARDUS. It holds the user's coin collection in encrypted local storage, orchestrates the off-chain protocols of \Cref{sec:mint-blindsig,sec:refresh} on behalf of the user, encodes payment requests as URIs, communicates with the encrypted relay mesh, and prepares on-chain transactions for the user's signature. This section specifies the wallet's data model, the invoice scheme, the backup and recovery procedures, the multi-device synchronisation discipline, and the hardware-wallet integration roadmap.

\subsection{Coin Storage}
\label{sec:wallet-coin-storage}

A wallet maintains a \emph{coin store}: an append-and-burn collection of $(x_i, Cp_i, \sigma_i, \mathsf{denom}_i)$ tuples, with a per-coin marker indicating whether the coin is spent (added to a nullifier on-chain), in-flight (currently the subject of a pending refresh or withdrawal), or active (spendable).

The store is held in encrypted local storage at all times. The encryption key is derived from the user's master secret (\Cref{sec:wallet-backup}). The encryption primitive is \texttt{ChaCha20-Poly1305} AEAD~\cite{nist-sp800-185}, with a 32-byte key, a 12-byte nonce per ciphertext, and an additional-data field carrying the wallet's protocol version and the user-chosen wallet identifier.

\subsection{Encrypted Backup and Recovery}
\label{sec:wallet-backup}

\paragraph{Backup format.} A backup is a single sealed blob:
\[
  \mathsf{Backup} = \mathsf{AEAD\text{-}Seal}(\mathsf{key} = \mathsf{HKDF}(\mathsf{master\_seed}, \text{``TARDUS-wallet-backup-v1''}), \mathsf{nonce}, \mathsf{aad}, \mathsf{plaintext})
\]
where $\mathsf{plaintext}$ is the borsh-encoded serialised coin store plus wallet metadata (version, creation timestamp, user-chosen label). The $\mathsf{master\_seed}$ is derived from a BIP-39 mnemonic (\Cref{sec:wallet-mnemonic}).

\paragraph{Persistent file format (v3.4).} The wallet ships with a
\texttt{WalletDb} struct that holds the coin store on disk at
\texttt{<data\_dir>/wallet.bin}. The wire format is the same
\texttt{Backup} blob above; the only addition is atomic write
discipline (temp file + rename) so a crash mid-write cannot leave a
half-truncated wallet file. The persistent format includes
\texttt{CoinStatus} (Active / InFlight / Spent) per stored coin, used
by the wallet UI for balance display and the refresh flow.

\paragraph{Multi-keyset metadata (v3.6).} A second AEAD file
\texttt{<data\_dir>/keysets.bin} holds a
\texttt{BTreeMap<\text{name}, KeysetInfo>} where each entry stores
\texttt{joint\_pk\_hex}, \texttt{denom},
\texttt{validators: Vec<URL>}, and optional CA / client cert paths
for TLS / mTLS communication. The same wallet seed encrypts both
files; the wallet UI lets the user issue or refresh against a named
keyset (\texttt{tardus-wallet issue --keyset mainnet}) without
re-typing validator URLs or joint pubkeys.

\subsection{Mnemonic Key Management}
\label{sec:wallet-mnemonic}

The wallet's $\mathsf{master\_seed}$ is derived from a BIP-39
mnemonic phrase~\cite{bip39}:

\begin{enumerate}[leftmargin=2em,itemsep=0pt]
  \item User generates a fresh 12- or 24-word phrase from
        $\mathsf{OsRng}$.
  \item PBKDF2-HMAC-SHA-512(mnemonic, optional passphrase, 2048 iterations)
        outputs 64 bytes.
  \item The wallet takes the \emph{first 32 bytes} as
        $\mathsf{master\_seed}$. The remaining 32 bytes are intentionally
        discarded; the wallet does not use BIP-32 hardened derivation
        paths because each coin's per-coin secret is derived via
        per-coin HKDF (\Cref{sec:refresh-derivation}) rather than via
        a deterministic tree.
\end{enumerate}

The BIP-39 test vector \texttt{"abandon abandon ... about"} with no
passphrase yields a \texttt{master\_seed} prefix
\texttt{5eb00bbddcf069084889a8ab9155568165f5c453ccb85e70811aaed6f6da5fc1};
the wallet's unit tests pin this vector for regression protection.

\subsection{Receiving Identity}
\label{sec:wallet-recv-identity}

For peer-to-peer payment delivery (\Cref{sec:relay}), a recipient
publishes a long-lived \emph{receiving identity} public key,
deterministically derived from $\mathsf{master\_seed}$:
\[
  \mathsf{recv\_sk} = \mathsf{HKDF\text{-}SHA\text{-}512}(\,\mathsf{salt}=\text{``TARDUS-recv-id-v1''},\,
                                                          \mathsf{ikm}=\mathsf{master\_seed},\,
                                                          \mathsf{info}=\text{``ed25519-keypair''}\,)\bmod \ell,
\]
$\mathsf{recv\_pk} = \mathsf{recv\_sk} \cdot G$. The pubkey is shared
with senders via the \texttt{recipient\_pubkey} field of the invoice
URI; the senders encrypt their coin-transfer payload to
$\mathsf{recv\_pk}$ via the sealed-box construction of
\Cref{sec:relay-payload}. The wallet never re-derives or rotates
$\mathsf{recv\_sk}$; recovery is full as long as the mnemonic is
available.

\paragraph{Why deterministic.} Recipients are typically off-line at the
moment of receipt; the relay holds the ciphertext until the recipient
polls. If the recipient's wallet could not re-derive
$\mathsf{recv\_sk}$ from the mnemonic alone, lost wallet devices would
lose every coin in flight at the moment of loss. The HKDF derivation
above costs $O(1)$ at boot and is bit-equivalent across devices.

\paragraph{Why separated from per-coin secrets.} The receiving
identity is a long-lived address; per-coin secrets are one-shot. Mixing
them via a single derivation tree (e.g.\ BIP-32) would let a sender who
sees one coin secret retroactively de-anonymise the recipient's other
addresses. The receiving identity therefore lives in its own
single-purpose HKDF context (salt = \texttt{"TARDUS-recv-id-v1"}),
disjoint from coin-secret derivation.

\paragraph{Recovery.} On a fresh device, the user enters the BIP-39 mnemonic; the wallet reconstructs $\mathsf{master\_seed}$, retrieves the latest backup (from the cloud or local copy), and decrypts the coin store. Coins recovered this way are unchanged: they have not been spent on-chain (since the coin store has not been re-used), so they retain their full value. \emph{Caveat:} backups are point-in-time snapshots. Coins minted or refreshed between the last backup and a device loss are unrecoverable. The wallet MUST trigger an incremental backup after every refresh.

\paragraph{Coin recovery is NOT deterministic regeneration.} Per the Cashu NUT-13 lesson (\texttt{research/PRODUCTION\_LESSONS.md} §L2), coin secrets are derived from per-coin one-shot HKDF seeds (\Cref{sec:refresh-derivation}), not from the wallet's master seed via BIP-32. A lost wallet whose backup is also lost is permanently unrecoverable; this trade-off is the price of soundness against keyset-collision attacks.

\subsection{Invoice URI Scheme}
\label{sec:wallet-invoice}

The TARDUS invoice URI captures all data needed for a payer to send funds to a recipient. The scheme is:
\[
  \texttt{tardus://}\langle\mathit{recipient\_pubkey}\rangle\texttt{?denom=}\langle d\rangle\texttt{\&relay=}\langle\mathit{relay\_url}\rangle\texttt{\&memo=}\langle\mathit{base64}\rangle
\]
with:
\begin{itemize}[leftmargin=2em,itemsep=0pt]
  \item $\mathit{recipient\_pubkey}$: $32$-byte recipient identity, hex-encoded.
  \item $\mathit{denom}$: requested amount in lamports.
  \item $\mathit{relay\_url}$: URL of an encrypted relay (\Cref{sec:wallet-relay}) the recipient is listening on. May appear multiple times.
  \item $\mathit{memo}$ (optional): base64-encoded free-form message, $\leq 128$ bytes after decoding.
\end{itemize}

The URI is presentable as a QR code, hyperlink, or NFC tag. Parsers MUST reject malformed URIs (wrong scheme, missing required parameters, exceeding length bounds).

\subsection{Refresh Session Persistence}
\label{sec:wallet-refresh-persistence}

The user-side state across the six-round refresh protocol (\Cref{sec:refresh-rounds}) is held in a \texttt{RefreshSession} record. Between Round 2 (submission of $\{c_\gamma\}$) and Round 6 (unblinding), the wallet MUST persist this record to encrypted local storage on every state transition, so a device crash between rounds does not strand the user with a coin marked in-flight but no recovery path.

The record is borsh-encoded and AEAD-sealed with the same key as the coin store. After Round 6 succeeds (the new coin is constructed), the session record is deleted and the new coin is committed to the coin store.

\subsection{Multi-Device Synchronisation}
\label{sec:wallet-multi-device}

Multi-device support is deferred to v1.5.2. The intended design uses a CRDT-based replicated coin store:

\begin{itemize}[leftmargin=2em,itemsep=0pt]
  \item Coin additions: G-Set of $Cp$ entries (monotonic merge).
  \item Coin spends: LWW-Register per $Cp$ with timestamp = epoch + slot of the spend transaction.
\end{itemize}

Cross-device synchronisation rides over the same encrypted relay mesh used for invoice payment delivery, with each device having its own relay handle.

\subsection{Relay Communication}
\label{sec:wallet-relay}

The wallet's user-to-user delivery path runs through the relay layer
specified in full in \Cref{sec:relay}. From the wallet's perspective,
the protocol is two commands:

\begin{enumerate}[leftmargin=2em,itemsep=0pt]
  \item \textbf{Send.} \texttt{tardus-wallet pay --invoice <uri> --keyset
        <name> [--encrypt]}: the wallet mints a fresh coin under the
        named keyset, JSON-encodes the coin material, optionally seals
        it to the invoice's recipient pubkey via the
        \Cref{sec:relay-payload} sealed-box construction, and
        \texttt{POST}s the resulting hex payload to the relay
        identified in the invoice URI.
  \item \textbf{Receive.} \texttt{tardus-wallet receive --relay <url>
        --recipient-pubkey <hex>}: the wallet polls the relay for
        messages addressed to the recipient pubkey, tries to
        \texttt{sealed\_box::open} each payload with the
        mnemonic-derived $\mathsf{recv\_sk}$ (with a graceful
        fallback to plain JSON for v5.3 senders), reconstructs the
        coin, adds it to the coin store, and (unless
        \texttt{--keep-on-relay} is set) deletes the message from the
        relay.
\end{enumerate}

Multiple relays may be used in parallel for redundancy; the wallet
polls each independently. The relay's TLS, persistence, and failure
semantics are specified in \Cref{sec:relay}.

\subsection{Hardware-Wallet Roadmap}
\label{sec:wallet-hw}

\begin{itemize}[leftmargin=2em,itemsep=0pt]
  \item \textbf{Phase 3 shipped (v3 wallet):} software-only wallets.
        The mnemonic is held by the user (offline backup); the wallet
        process holds the derived master seed only during the
        function call that needs it (\texttt{Zeroizing} wrapper, dropped
        at scope exit). \texttt{wallet.bin} + \texttt{keysets.bin} are
        AEAD-sealed under HKDF-derived keys.
  \item \textbf{Phase 4 (planned):} Tangem and Trezor support via
        custom apps (blind Schnorr signing implemented on the hardware
        device). YubiKey 5 with FIDO2 for Ed25519 signing operations
        on the receiving identity.
  \item \textbf{Phase 5 (planned):} Ledger Nano support, contingent on
        Ledger firmware acceptance of TARDUS blind-Schnorr signing.
        Sealed-box decryption may remain on the host (X25519 not
        always available on HW wallets); the recipient's secret key
        is then derived per-call and zeroised, accepting a smaller
        attack surface than long-lived in-process key residency.
\end{itemize}

\subsection{Security Considerations}
\label{sec:wallet-security}

The wallet is the user's trust anchor. Compromise of the wallet master seed compromises all coins held in that wallet. The wallet MUST:

\begin{itemize}[leftmargin=2em,itemsep=0pt]
  \item Never log, print, or otherwise expose secret material (coin secrets, refresh session blinding factors, master seed).
  \item Verify every received coin against the on-chain keyset registry's current $\mathsf{joint\_pk}$ for the claimed denomination before accepting it into the coin store.
  \item Detect and reject double-receipt: the same coin received twice (via two different relays, say) must merge to a single entry rather than appear twice.
  \item Enforce per-coin spend marking atomically across persistence operations.
\end{itemize}

\section{Security Model and Theorems}
\label{sec:security}

This section formalises the adversary model, lists the cryptographic and
system assumptions on which TARDUS rests, and states the six main theorems
together with reduction sketches. Full proof mechanisation is deferred to
Phase 4 (EasyCrypt); the proof sketches given here are sufficient to guide
that work and to serve as the audit baseline.

\subsection{Threat Model}
\label{sec:threat-model}

\paragraph{Adversary classes.} TARDUS considers four overlapping adversary
classes; each theorem in this section specifies which subset it tolerates.

\begin{description}[leftmargin=2em,style=nextline]
  \item[$\Adv_{\mathsf{net}}$ (passive observer)]
    Sees every on-chain transaction, all relay traffic between the user and
    the mint committee, and the contents of the public registries
    (\S5). Cannot inject, drop, or alter messages. This is
    the strictly weakest model.
  \item[$\Adv_{\mathsf{user}}$ (malicious user)]
    A polynomial-time party that controls one or more user identities, can
    interact with the mint committee through the canonical protocols
    (\S3 and \S4), and can submit arbitrary
    transactions to the on-chain program. Adaptive: may corrupt mint
    operators after seeing earlier transcripts.
  \item[$\Adv_{\mathsf{mint},t-1}$ (corrupted mint minority)]
    Statically corrupts up to $t-1$ of the $n$ mint operators (where $t$ is
    the signing threshold). Sees all corrupted operators' shares, internal
    state, and HSM-mediated traffic. Honest operators communicate over
    authenticated channels with replay protection. This is the maximum
    corruption tolerated by all positive results below; at $t$ corruptions
    coin forgery becomes possible by reconstruction.
  \item[$\Adv_{\mathsf{chain}}$ (network-level adversary)]
    Controls the Solana network up to but not including consensus
    (cf.\ Solana's $1/3$ honest-stake assumption). Can reorder
    transactions within a slot, drop unsubmitted transactions, and observe
    all account contents. Cannot subvert the on-chain program's atomic
    state transitions once a transaction is included in a finalised block.
\end{description}

\paragraph{Composability.} Privacy results in this section hold against the
union $\Adv_{\mathsf{net}} \cup \Adv_{\mathsf{mint},t-1} \cup
\Adv_{\mathsf{chain}}$ jointly; double-spend resistance holds additionally
against $\Adv_{\mathsf{user}}$.

\paragraph{What is \emph{not} modelled.} We exclude side-channel attacks on
the user's wallet device (covered by wallet-level mitigations in
\S6), targeted denial-of-service against individual mint
operators (covered by the $t$-out-of-$n$ availability margin in
\S8), and physical compromise of HSMs (covered by the
mandatory FIPS 140-2 Level 3 deployment requirement in \S3.1).

\subsection{Security Assumptions}
\label{sec:security-assumptions}

\paragraph{A1 (ECDLP on $\mathbb{G}$).}
\label{assn:ecdlp}
For the ed25519 prime-order subgroup $\mathbb{G}$ of order
$\ell \approx 2^{252}$, given a uniformly random $X = x \cdot G$, no
polynomial-time adversary recovers $x$ with non-negligible probability:
\[
  \Pr_{x \gets \mathbb{Z}_\ell}\bigl[\Adv(X) = x\bigr] \leq \negl(\lambda).
\]

\paragraph{A2 (Random-Oracle Model).}
\label{assn:rom}
We model SHA-512-reduced-mod-$\ell$ (used for Schnorr challenges,
\S2.3) and SHA-256 (used for nullifiers,
\S4.2) as random oracles in the security proofs of
T1, T2, T4. The hash-to-scalar output bias is bounded by $2^{-260}$
(\S2.3), which is sub-negligible for any
$\poly(\lambda)$ adversary.

\paragraph{A3 (Pedersen binding under DDH on $(G, H)$).}
\label{assn:pedersen}
For the joint commitment $C_{i,0} = s_i \cdot G + r_i \cdot H$
(\S3.5) with independently sampled generator $H$
(canonical hash-to-curve of the domain string \texttt{"TARDUS-vss-h-v1"}),
no polynomial-time adversary opens $C$ to two distinct pairs
$(s_i, r_i), (s_i', r_i')$ with $s_i \neq s_i'$.

\paragraph{A4 (Hash collision resistance).}
\label{assn:cr}
SHA-256 and SHA-512 are collision-resistant against $\poly(\lambda)$
adversaries with output-bit-length security.

\paragraph{A5 (Honest threshold).}
\label{assn:threshold}
In any DKG, sign, or refresh ceremony, at least $t$ of the $n$ mint
operators follow the protocol honestly. Corruption is static
(decided before the ceremony begins); the proactive reshare protocol
(\S3.7) re-randomises shares between epochs to mitigate
slow-creep corruption.

\paragraph{A6 (Solana consensus safety).}
\label{assn:solana}
Transactions included in a Solana block finalised at depth $\geq 32$ slots
are durable; the on-chain program's atomic state transitions hold under
this finality assumption.

\subsection{Theorem T1 --- Coin Unforgeability}
\label{sec:thm-unforgeability}

\begin{theorem}[Coin Unforgeability]
\label{thm:unforgeability}
Under assumptions A1 (ECDLP), A2 (ROM), and A5 (honest threshold), the
TARDUS threshold blind Schnorr signature (\S4) is
existentially unforgeable under chosen-message attack
($\mathsf{EUF\text{-}CMA}$) against any adversary $\Adv \in
\Adv_{\mathsf{user}} \cup \Adv_{\mathsf{mint},t-1}$. Formally, for any
$\poly(\lambda)$-time $\Adv$ making $q_s$ signing queries and $q_h$
random-oracle queries,
\[
  \Pr\bigl[\AdvG{\mathsf{unforge}}^{\Adv}(\lambda) = 1\bigr] \leq
  \frac{q_h \cdot \mathsf{Adv}^{\mathsf{ECDLP}}_{\mathbb{G}}(\lambda)}{1 - 1/\ell}
  + \negl(\lambda).
\]
\end{theorem}

\paragraph{Reduction sketch.} The proof composes two classical results:

\begin{enumerate}[leftmargin=2em,itemsep=0pt]
  \item \emph{Single-key Schnorr is EUF-CMA under ECDLP + ROM} by
        Pointcheval-Stern (\cite{pointcheval-stern1996}) via the
        forking-lemma rewinding argument.
  \item \emph{Threshold lift preserves unforgeability} by Stinson-Strobl
        (\cite{stinson-strobl2001}) when the underlying blind signature is
        unforgeable and the threshold is at most $t-1$. The reduction
        simulates the corrupted minority's view from the random
        oracle programming, exactly as in single-key.
\end{enumerate}

The validator-side commitment phase
(\S4.4) prevents the bias attack of
\cite{nick-ruffing-seurin2020}: validators publish $R_i = k_i \cdot G$
\emph{before} seeing the user's message, so the adversary cannot adaptively
choose challenges to favour a target signature.

\paragraph{Corruption boundary.} At exactly $t$ corruptions the adversary
reconstructs the joint secret $x = \sum_{i \in S} \lambda_{S,i} \cdot s_i$
via Lagrange interpolation (\S3.6) and forgery is trivial.
The proactive reshare protocol (\S3.7) bounds the corruption
window to one epoch.

\subsection{Theorem T2 --- Issuance Blindness}
\label{sec:thm-blindness}

\begin{theorem}[Issuance Blindness]
\label{thm:issuance-blindness}
Under assumption A1 (ECDLP) and the Pointcheval-Stern blinding analysis
(\cite{pointcheval-stern1996}), the user's view in a TARDUS issuance
session is statistically independent of the resulting coin's pair
$(Cp, \sigma)$. Formally, for any unbounded mint-side adversary
$\Adv \in \Adv_{\mathsf{mint},t-1}$ and any two messages
$Cp_0, Cp_1 \in \mathbb{G}$, the issuance transcripts produced by signing
$Cp_0$ and $Cp_1$ are perfectly indistinguishable:
\[
  \mathsf{SD}\bigl(\mathsf{Trans}(Cp_0),\ \mathsf{Trans}(Cp_1)\bigr) = 0.
\]
\end{theorem}

\paragraph{Reduction sketch.} The user samples uniformly random blinding
factors $\alpha, \beta \gets \mathbb{Z}_\ell^\ast$ in Round 2 of the
threshold blind sign (\S4.5). The blinded challenge
$c' = c + \beta$ and blinded commitment $R' = R + \alpha \cdot G + \beta
\cdot Cp$ are uniform over $\mathbb{Z}_\ell$ and $\mathbb{G}$
respectively, independent of $Cp$. The mint sees only $(R', c', s')$
where $s'$ is the response to $c'$; the unblinding step
($s = s' - \alpha$, $c$ recovered) preserves uniformity. The
information-theoretic argument gives statistical distance $0$, stronger
than the computational indistinguishability typically achieved by
zero-knowledge schemes.

\paragraph{Note.} The blindness result is \emph{unconditional} (no
computational assumption on the mint side), in contrast to T1 which
reduces to ECDLP. This is a strength inherited from Chaum's 1982
construction.

\subsection{Theorem T3 --- Double-Spend Prevention}
\label{sec:thm-double-spend}

\begin{theorem}[Double-Spend Prevention]
\label{thm:double-spend}
Under assumptions A4 (SHA-256 collision resistance) and A6 (Solana
finality), any attempt by $\Adv \in \Adv_{\mathsf{user}}$ to submit a
spend for the same coin $(Cp, \sigma)$ twice is rejected with probability
$1$ at the protocol level. Formally,
\[
  \Pr\bigl[\AdvG{\mathsf{double\text{-}spend}}^{\Adv}(\lambda) = 1\bigr] = 0
\]
after the first spend's transaction has finalised on chain.
\end{theorem}

\paragraph{Proof.} The nullifier
$\mathsf{null}(Cp) = \mathsf{SHA\text{-}256}(\text{``TARDUS-nullifier-v1''} \,\|\, Cp)$
is deterministic in $Cp$ (\S4.2). The on-chain
\texttt{Refresh} and \texttt{Withdraw} handlers (\S5.3.3)
perform $\mathsf{NullifierSet.insert}(\mathsf{null}(Cp))$, which returns
\texttt{false} if the entry is already present and aborts the
transaction with \texttt{ERR\_DOUBLE\_SPEND}. This insertion is wrapped
inside the atomic boundary of a single Solana transaction; under A6 the
state mutation is durable upon finalisation. Hash collisions are excluded
by A4. \qed

\paragraph{Empirical anchor.} The on-chain program v1.4.10 deployed at
\texttt{AmY1ysgQyCC6CmorXkrNkogBSHmomy4kbsPziAYUx47u} on Solana devnet
implements this exact check; the CLI's client-side preflight
(the CLI \texttt{tardus devnet refresh}) rejects double-spend attempts before paying
transaction fees, providing defence in depth.

\subsection{Theorem T4 --- Cut-and-Choose Soundness}
\label{sec:thm-cut-and-choose}

\begin{theorem}[Cut-and-Choose Soundness for Refresh]
\label{thm:cut-and-choose}
Under assumptions A1 (ECDLP) and A2 (ROM), the $\kappa$-fold
cut-and-choose refresh protocol (\S4) binds a malicious
user to producing a denomination-preserving and unlinkable replacement coin
with cheating probability at most
\[
  \Pr[\mathsf{cheat}] \leq \frac{1}{\kappa + 1}.
\]
For the default $\kappa = 32$ (\S4.5), this is
$\approx 3.0\%$ per session; a malicious user wishing to extract
$\Delta$-fold value must succeed $\Delta$ times, with probability
$1 / 33^\Delta$.
\end{theorem}

\paragraph{Proof sketch.} The user commits in Round 2 to $\kappa + 1$
candidate trapdoors $(c_1, \ldots, c_{\kappa+1})$, each encoding a
candidate replacement coin. The mint's Round 3 random challenge selects a
subset $S \subset [\kappa+1]$ of size $\kappa$ to reveal. The remaining
candidate $j = [\kappa+1] \setminus S$ is unblinded and signed; the user
recovers a single fresh coin. For the user to cheat (e.g.\ extract more
denominations than they surrendered), all $\kappa$ revealed candidates
must be consistent with the protocol \emph{and} the one unrevealed
candidate must encode the cheat: probability $1/(\kappa+1)$ assuming the
mint's challenge is uniform. The Round 4 reveal-and-verify step
(\S4.5) checks consistency cryptographically; the
Pedersen binding under A3 prevents equivocation between commit and reveal.

\paragraph{Anonymity preservation.} The unrevealed candidate $j$ is
information-theoretically independent of the user's identity (it is a
freshly sampled scalar). The mint sees only the revealed $\kappa$
candidates, which the user has bound by commitment, plus the blinded
issuance transcript of $j$ which by T2 is independent of the underlying
$(Cp_{\mathsf{new}}, \sigma_{\mathsf{new}})$. Hence the new coin is
unlinkable to the old coin in the union of $\Adv_{\mathsf{net}}$ and
$\Adv_{\mathsf{mint},t-1}$.

\subsection{Theorem T5 --- Reshare Correctness}
\label{sec:thm-reshare}

\begin{theorem}[Proactive Reshare Joint-Key Invariance]
\label{thm:reshare}
Under assumptions A1 (ECDLP), A3 (Pedersen binding), and A5 (honest
threshold), the proactive reshare protocol (\S3.7) transforms
the shares $(s_1, \ldots, s_n)$ into fresh shares $(s_1', \ldots, s_n')$
such that:
\begin{enumerate}[leftmargin=2em,itemsep=0pt]
  \item the joint public key $X = x \cdot G$ is preserved (where
        $x = \sum_{i} \lambda_i s_i = \sum_{i} \lambda_i s_i'$);
  \item the new shares are uniformly distributed conditional on $X$;
  \item old shares from epoch $e$ are useless for any operation in epoch
        $e + 1$ or later.
\end{enumerate}
\end{theorem}

\paragraph{Proof sketch.} The reshare uses a zero-polynomial Pedersen-VSS:
each operator distributes shares of $0$ via a fresh degree-$(t-1)$
polynomial $\tilde{f}_i$ with $\tilde{f}_i(0) = 0$. New shares are
$s_j' = s_j + \sum_i \tilde{f}_i(j)$. The dual-commitment broadcast
(\S3.5) lets honest operators detect a cheating
dealer who chose $\tilde{f}_i(0) \neq 0$; the additional check
$\tilde{A}_{i,0} = \mathcal{O}$ (the identity element)
catches this case explicitly (\S3.7). The
joint-key invariance is algebraic: $\sum_j \lambda_j s_j' =
\sum_j \lambda_j s_j + \sum_i \tilde{f}_i(0) = x + 0 = x$. The
uniformity follows from the freshness of $\tilde{f}_i$. The
\emph{usefulness} of old shares for epoch $e+1$ requires solving
$s_j' = s_j + \delta_j$ for $\delta_j$ given only $X$; this is equivalent
to ECDLP. \qed

\paragraph{Empirical anchor.} The crown-jewel
\texttt{sign\_after\_reshare} test in the \texttt{tardus-mint} crate
(test suite) demonstrates that a signature produced by a
$t$-subset of the new shares verifies under the original $X$, while a
signature attempted with one new and one old share fails.

\subsection{Theorem T6 --- Vault Collateral Invariant}
\label{sec:thm-collateral}

\begin{theorem}[Vault Collateral Invariant]
\label{thm:collateral}
Under assumption A6 (Solana finality) and the on-chain handlers'
canonical-PDA validation (\S5.4), at every
finalised state $\sigma$ of the program the vault's lamport balance
satisfies the bookkeeping invariant
\[
  \mathsf{Vault}[\Denom].\mathsf{lamports} \;\geq\;
  \Denom \cdot \bigl(\,
    \#\mathsf{IssuedCoins}[\Denom] - \#\mathsf{NullifierSet} \cap
        \mathsf{NullifiersOf}[\Denom]
  \,\bigr).
\]
\end{theorem}

\paragraph{Proof.} The handler-by-handler accounting is mechanical:
\texttt{Deposit} requires a preceding system\_program::transfer of at least
$\Denom$ lamports to the vault PDA, observed via
\texttt{vault\_account.lamports()} (\S5.3.2);
\texttt{Withdraw} removes exactly $\Denom$ lamports via
\texttt{invoke\_signed} after a successful nullifier insertion
(\S5.3.4); \texttt{Refresh} does not touch the vault
(\S5.3.3, anonymity preservation). The on-chain code
explicitly checks
\texttt{vault\_account.lamports() < denom} and returns
\texttt{ERR\_VAULT\_INSUFFICIENT} if violated. Hence the invariant is
restored after every successful transaction; A6 ensures durability. \qed

\paragraph{Operational implication.} A user can perform a
\emph{collateral audit} at any time by reading the vault account's
lamport balance and the on-chain nullifier set, then comparing against
the off-chain issued-coin database (publicly aggregated by the mint
committee as part of the transparency log). Any deviation is publicly
observable.

\subsection{Theorem T7 --- Sealed-Box Payload Confidentiality}
\label{sec:theorem-T7}

\begin{theorem}[Sealed-Box IND-CCA, v1.7]
\label{thm:T7}
Let $\mathsf{SB}$ be the TARDUS sealed-box AEAD construction of
\Cref{sec:relay-payload}: $\mathsf{Seal}(m, pk_R) = epk \,\|\,
\mathsf{ChaCha20\text{-}Poly1305}_{K,N}(m)$ where
$(K,N) = \mathsf{HKDF\text{-}SHA\text{-}256}(\mathsf{ECDH}(esk, pk_R^{(\mathrm{X25519})}))$.
For any PPT adversary $\mathcal{A}$ playing the IND-CCA game and
controlling the relay operator (i.e. observing all sealed-box wire
bytes and adaptively requesting decryptions of any non-challenge
ciphertext under the recipient's secret key), there exist PPT
adversaries $\mathcal{B}_1, \mathcal{B}_2$ such that
\[
  \mathsf{Adv}^{\mathrm{IND\text{-}CCA}}_{\mathsf{SB}}(\mathcal{A}) \;\leq\;
  \mathsf{Adv}^{\mathrm{IND\text{-}CCA}}_{\mathsf{ChaCha20\text{-}Poly1305}}(\mathcal{B}_1) \;+\;
  \mathsf{Adv}^{\mathrm{ECDLP}}_{X25519}(\mathcal{B}_2) \;+\;
  \frac{q_H}{2^{256}} \;+\; \mathsf{negl}(\lambda),
\]
where $q_H$ is the number of HKDF (i.e.\ SHA-256-based PRF)
queries.
\end{theorem}

\begin{proof}[Proof sketch]
The reduction composes three layers, each modular against the
others:

\begin{enumerate}[leftmargin=2em,itemsep=0pt]
  \item \textbf{ECDH key indistinguishability.}
    Replace the genuine shared secret
    $s = \mathsf{ECDH}(esk, pk_R^{(\mathrm{X25519})}) = esk \cdot pk_R^{(\mathrm{X25519})}$
    with a uniformly random
    $s' \xleftarrow{\$} \{0,1\}^{256}$ (the ECDLP / DDH game on
    Curve25519). An adversary distinguishing genuine from random
    $s$ in the sealed-box context wins the DDH game for X25519;
    the loss is $\mathsf{Adv}^{\mathrm{ECDLP}}_{X25519}(\mathcal{B}_2)$.
  \item \textbf{HKDF PRF.}
    With $s'$ uniform, the HKDF-derived $(K, N)$ pair is
    statistically indistinguishable from uniform up to a
    $q_H / 2^{256}$ collision probability across the
    $\mathcal{A}$'s HKDF queries.
  \item \textbf{ChaCha20-Poly1305 IND-CCA.}
    The challenge ciphertext is then a fresh ChaCha20-Poly1305
    encryption under uniformly random $(K, N)$. Any IND-CCA
    distinguisher on the sealed-box reduces directly to an IND-CCA
    distinguisher on the underlying AEAD; the loss is
    $\mathsf{Adv}^{\mathrm{IND\text{-}CCA}}_{\mathsf{ChaCha20\text{-}Poly1305}}(\mathcal{B}_1)$.
\end{enumerate}

The compromised-relay adversary model is captured by giving
$\mathcal{A}$ the oracle that returns sealed-box wire bytes for
arbitrary plaintexts (chosen-plaintext oracle) plus the
decryption oracle for any non-challenge ciphertext (chosen-
ciphertext oracle). Both are simulated against the underlying
ChaCha20-Poly1305 oracle in step 3 after the ECDH-to-random hop.

The full mechanisation is in
\texttt{spec/easycrypt/sealed\_box.ec} (Phase 4 audit firm
hand-off).
\end{proof}

\subsection{Theorem T8 --- Confidential Vault Indistinguishability (v1.7 placeholder)}
\label{sec:theorem-T8}

\begin{theorem}[Vault Balance Indistinguishability, v1.5+ target]
\label{thm:T8}
Let $\mathsf{Vault}^{\mathrm{conf}}$ be the v1.5+
Token-2022 ConfidentialMint vault construction (deferred; design
doc at \texttt{deploy/runbooks/token-2022-confidential-mint-design.md}).
For any PPT adversary $\mathcal{A}$ observing the on-chain
encrypted balance commitments + auditor public key + all public
TARDUS metadata, the probability $\mathcal{A}$ distinguishes the
real vault SOL balance $a_0$ from any
$a_1 \in [0, 2^{64})$ is bounded by:
\[
  \mathsf{Adv}^{\mathrm{IND}}_{\mathsf{Vault}^{\mathrm{conf}}}(\mathcal{A}) \;\leq\;
  \mathsf{Adv}^{\mathrm{IND\text{-}CPA}}_{\mathsf{ElGamal}_{\mathrm{Edwards}}}(\mathcal{B}_3) \;+\;
  \mathsf{Adv}^{\mathrm{Soundness}}_{\mathsf{RangeProof}}(\mathcal{B}_4) \;+\;
  \mathsf{negl}(\lambda).
\]
\end{theorem}

\paragraph{Status.} \textbf{Placeholder; not provable on v1
plain-SOL vault.} v1 ships with publicly visible vault SOL
balances (the deposit and withdraw amounts are observable on
the Solana chain). T8 becomes provable once Faz E full ships
(estimated v1.5+, $\approx 19$ engineering weeks; see
architecture decision record in
\texttt{token-2022-confidential-mint-design.md}). The placeholder
is included here to lock the audit firm's expected formal-
verification surface area pre-engagement, so the Phase 4 audit
contract can reference T8 by name.

\paragraph{Reduction sketch (for v1.5+).}
The reduction composes Token-2022's
ElGamal-encrypted-balance ciphertext (the $\mathcal{B}_3$ hop)
with the deposit / withdraw range-proof soundness
($\mathcal{B}_4$ hop). Concretely: the auditor's view sees
$\mathsf{Enc}_{\mathrm{auditor}}(a)$; an IND-CPA distinguisher on
ElGamal under the Edwards subgroup recovers $a$ vs $a'$ with the
$\mathcal{B}_3$ advantage. The range proof's soundness ensures
the amount is in the declared range $[0, 2^{64})$; without
soundness, an adversary could mint negative amounts and break
T6 simultaneously.

\subsection{Theorem T9 --- Fee-Payer Unlinkability (v1.8)}
\label{sec:theorem-T9}

\begin{theorem}[Fee-Payer Unlinkability, Layer 1--4 composed]
\label{thm:T9}
Let $\mathcal{P} = \{P_1, \ldots, P_n\}$ be the set of TARDUS
spend-class transactions (Refresh or Withdraw) submitted to
Solana within an observation window $W$, each accompanied by a
funding chain through the on-chain $\mathsf{SponsorPool}$ PDA per
\Cref{sec:mint-fee-payer-privacy} (Layer 3+4). Let
$\mathcal{F} = \{F_1, \ldots, F_m\}$ be the set of distinct
deposit-class transactions to the $\mathsf{SponsorPool}$ inside
the same window $W$, originating from independent funding
wallets.

For any PPT chain-graph adversary $\mathcal{A}$ with read access
to the entire Solana ledger state inside $W$ — including every
TX hash, every account balance trajectory, every instruction
log, and the public IP source of each TX submission — the
advantage of $\mathcal{A}$ in correctly mapping any specific
spend-class TX $P_i$ to the funding wallet whose deposit
ultimately funded its fee payment is bounded by:
\[
  \mathsf{Adv}^{\mathrm{link}}_{\mathsf{Layer1{+}4}}(\mathcal{A}) \;\leq\;
  \frac{1}{|\mathcal{F}|} \;+\;
  \mathsf{Adv}^{\mathrm{IP\text{-}deanon}}_{\mathrm{network}}(\mathcal{B}_5) \;+\;
  \mathsf{Adv}^{\mathrm{timing}}_{\mathrm{slot}}(\mathcal{B}_6) \;+\;
  \mathsf{negl}(\lambda),
\]
where $\mathcal{B}_5$ wins an IP-deanonymisation game (e.g.\
linking submitter IPs across the Solana RPC layer) and
$\mathcal{B}_6$ exploits per-slot timing fingerprints across the
deposit and payout TXs.
\end{theorem}

\begin{proof}[Proof sketch]
The reduction has three layers:

\begin{enumerate}[leftmargin=2em,itemsep=0pt]
  \item \textbf{Pool indistinguishability ($1/|\mathcal{F}|$ term).}
    The $\mathsf{SponsorPool}$ PDA is a single program-owned account
    whose balance trajectory aggregates all deposits and all
    payouts. From the chain observer's view, every deposit
    $F_j$ contributes its lamports indistinguishably to the
    pool's running balance; every payout $\mathsf{SponsorPayout}$
    instruction (\Cref{sec:onchain-instr-sponsor}) drains a
    fixed amount $\Lambda$ to an ephemeral keypair $E_k$. Inside
    the on-chain log, $E_k$ has no traceable predecessor beyond
    "drawn from the pool"; the funding wallet $F_j$ that
    notionally "covered" the payout is statistically uniform
    over $\mathcal{F}$ as long as deposits arrive faster than
    payouts (a steady-state liquidity assumption).

  \item \textbf{Ephemeral signer + pool caller rotation (Layers 1, 2, 4).}
    Each spend-class TX $P_i$ is signed by a fresh ephemeral
    keypair $E_i$ (Faz 9.1) funded by the pool payout. Each
    pool payout itself is invoked by a caller keypair drawn at
    random from the sponsor pool (Faz 9.5). The on-chain
    signing-chain graph
    $E_i \leftarrow \mathsf{Payout}_{C_j} \leftarrow F_k$
    has fan-in $|\mathcal{F}|$ from every spend's perspective,
    so the joint posterior $\Pr[E_i \text{ funded by } F_k \mid
    \text{public chain state}]$ is $\approx 1/|\mathcal{F}|$.

  \item \textbf{Side channels ($\mathcal{B}_5$, $\mathcal{B}_6$).}
    The reduction does not close (i) network-layer
    deanonymisation (the submitter's IP visible to the Solana
    RPC endpoint and any intermediate router), nor (ii) per-slot
    timing fingerprints (a deposit at slot $S$ and a payout at
    slot $S+\Delta$ for small $\Delta$ creates a correlation
    even if no direct funding link is on-chain). These are
    explicitly enumerated in
    \Cref{sec:mint-fee-payer-privacy} ("Residual leaks not
    addressed by Layers 1--4") and require independent
    countermeasures (Tor / I2P for $\mathcal{B}_5$; random
    submission delays + Jito bundle ordering of
    \Cref{sec:mint-batch} for $\mathcal{B}_6$).
\end{enumerate}

The deposit-rate-exceeds-payout-rate assumption is operationally
necessary: if the pool drains faster than it refills, the
adversary can pinpoint the specific deposit that "must have"
funded a specific payout by lamport-balance accounting. The
spec's drain-rate ceiling
($\approx 14\,\mathrm{SOL/hour}$ at default parameters; see
\Cref{sec:onchain-instr-sponsor}) is calibrated so this
condition is satisfied under typical mainnet deposit cadence.
\end{proof}

\paragraph{Status.} v1.8 spec — theorem statement and informal
reduction shipped. Full mechanisation in EasyCrypt is a Phase 4
audit firm task (the $1/|\mathcal{F}|$ uniform-posterior bound
is straightforward but the steady-state liquidity assumption
requires a real-world model that the audit firm scopes).

\paragraph{Concrete empirical bound (devnet).} The devnet
$\mathsf{SponsorPool}$ PDA holds $\approx 49\,000\,000$ lamports
across $\approx 5$ historical deposits, with payouts of
$10^6$ lamports each. With $|\mathcal{F}| = 5$, an adversary's
direct link probability is $\approx 20\%$ — improving to $< 1\%$
once $|\mathcal{F}| \geq 100$ on mainnet under healthy deposit
cadence. The live demo of 2026-05-22 (TX
\texttt{3UVkfK5tt671FDJ\ldots}) exhibits this property
end-to-end: the spend-class Withdraw signer
\texttt{B3mFCD2qdedYd2nYuTcwZFZS1ghCs7DwKQRtgrvhWszq} is
demonstrably an ephemeral keypair distinct from the deployer
wallet, funded via an on-chain $\mathsf{SponsorPayout}$.

\subsection{Composition: TARDUS-Level Security}
\label{sec:composition}

\begin{corollary}[End-to-End Security]
\label{cor:end-to-end}
The TARDUS protocol satisfies the following composite guarantees:
\begin{enumerate}[leftmargin=2em,itemsep=0pt]
  \item \textbf{Coin authenticity} (T1 + T5): no party can produce a valid
        $(Cp, \sigma)$ without participating in a legitimate issuance or
        refresh session, even across reshare epochs;
  \item \textbf{Sender/receiver/amount privacy} (T2 + T4): on-chain
        observers cannot link a refreshed coin to its predecessor, the
        mint cannot link an issued coin to a refreshed predecessor, and
        the per-coin denomination is the only amount-related leak;
  \item \textbf{Atomic ledger integrity} (T3 + T6): the on-chain vault
        balance and nullifier set are consistent at every finalised
        slot, and double-spends are deterministically rejected.
\end{enumerate}
\end{corollary}

\paragraph{Remaining gaps.} The denomination is leaked at deposit and
withdraw time (the user surrenders exactly $\Denom$ lamports to the vault
PDA, observable on-chain). Phase 1.4.11 will replace the SOL vault with
a Solana Token-2022 Confidential Mint vault, which encrypts the deposit
and withdraw amounts; the per-coin denomination then becomes a private
field of the mint's bookkeeping, not visible on the chain. Phase 1.4.12
will migrate the Borsh-encoded nullifier set
(\S5.3.3) to a Light Protocol compressed Merkle
tree, raising the per-vault capacity from $\approx 250$ to $2^{32}$
nullifiers.

\subsection{Mechanisation Roadmap}
\label{sec:mech-roadmap}

Phase 4 of the project will mechanise theorems T1, T2, and T4 in
EasyCrypt (\cite{easycrypt2013}). Tactical sketch:

\begin{description}[leftmargin=2em,style=nextline]
  \item[\texttt{schnorr.ec}]
    Encodes the single-key Schnorr signature, the
    Pointcheval-Stern forking argument, and the ROM tape. Discharges T1
    in the single-key case.
  \item[\texttt{threshold\_lift.ec}]
    Encodes the simulator for the corrupted minority's view (per
    \cite{stinson-strobl2001}) and reduces threshold T1 to single-key T1.
  \item[\texttt{blind.ec}]
    Encodes the blinding-factor algebra of T2 and proves perfect
    indistinguishability (zero statistical distance) of the issuance
    transcripts.
  \item[\texttt{cut\_choose.ec}]
    Encodes the $\kappa$-fold cut-and-choose game and proves the
    $1/(\kappa+1)$ cheating bound of T4.
\end{description}

T3, T5, and T6 are operational invariants whose mechanisation reduces to
runtime checks already present in the reference implementation; they
require no separate proof effort beyond the test suite of
test suite.

\section{Failure Modes and Recovery}
\label{sec:failure-modes}

This section catalogues the failure modes TARDUS is designed to survive,
specifies the recovery procedure for each, and documents the residual
classes that operators must monitor for. The catalogue is the operational
counterpart of \S7's adversary model: \S7 says \emph{which adversaries
the protocol resists}, \S8 says \emph{what happens when something actually
goes wrong} and \emph{how to detect and recover}.

Failure modes are grouped by origin: validator-side (\S8.1),
user-side (\S8.2), protocol-level adversarial (\S8.3), and
network-level (\S8.4). Each entry lists \textit{Trigger},
\textit{User-facing consequence}, \textit{Recovery}, and
\textit{Detection signal}.

\subsection{Validator-Side Failure Modes}
\label{sec:failure-validator}

\paragraph{F1 --- Single validator crash.}
\begin{description}[leftmargin=2em,style=nextline,itemsep=0pt]
  \item[Trigger] One of the $n$ mint operators is offline (hardware
    fault, network drop, process restart).
  \item[Consequence] If the remaining online operators number $\geq t$,
    the mint continues to issue, refresh, and rotate normally; the
    offline operator simply does not contribute to active sessions.
    If online $< t$, all new issuance and refresh sessions stall;
    \emph{withdrawals} continue without disruption (on-chain only).
  \item[Recovery] The crashed operator restarts, loads its persistent
    share state from encrypted disk (\S3.1 HSM-mediated storage),
    and rejoins the next ceremony. No share regeneration required.
  \item[Detection] Operator heartbeat absent from the transparency log
    for $>$ 30 s. The validator daemon (Faz 2) publishes a liveness
    signal every 10 s; three consecutive misses raise an operator
    on-call alert.
\end{description}

\paragraph{F2 --- Validator network partition.}
\begin{description}[leftmargin=2em,style=nextline,itemsep=0pt]
  \item[Trigger] A network split isolates a subset $P \subsetneq [n]$
    of operators from the rest.
  \item[Consequence] If $|P| < t$ and $|[n] \setminus P| \geq t$,
    the majority side continues; minority side stalls.
    If both sides are $\geq t$, neither side should sign autonomously
    (a brain-split risk for the joint key). The validator daemon's
    quorum check enforces this: each side independently observes
    the other side's absence and pauses.
  \item[Consequence (rare)] If both sides are $\geq t$ and both ignore
    the quorum check (operator misconfiguration), both sides can issue
    valid signatures under the same joint key, creating fork-style
    double-issuance. This is operationally catastrophic but the
    protocol still rejects on-chain double-spend (T3) — the conflict
    surfaces at the first user who tries to spend two coins that
    happen to share a nullifier ancestry. Recovery requires a manual
    reconciliation against the transparency log.
  \item[Recovery] Partition heals → next ceremony auto-resumes. If
    forked-issuance occurred, manual epoch rotation via a key-revoke
    + re-register cycle (\S5.3) under committee consensus.
  \item[Detection] Cross-side keepalive failure → operator on-call;
    automated quorum-check pause is the first line of defence.
\end{description}

\paragraph{F3 --- HSM failure.}
\begin{description}[leftmargin=2em,style=nextline,itemsep=0pt]
  \item[Trigger] An operator's FIPS 140-2 Level 3 HSM (\S3.1) hardware
    fails or the key wrap protocol detects tampering.
  \item[Consequence] That operator can no longer participate; treat as
    F1 (single-validator-crash) extended.
  \item[Recovery] Replace the HSM, re-initialise with a fresh share
    derived from a proactive reshare (\S3.7). The reshare protocol's
    zero-polynomial property (T5) means the joint key is preserved
    and the failed operator's old share becomes useless — even if
    later recovered from the failed HSM. Operator's share is replaced
    in-epoch; no user-facing impact beyond F1-class downtime.
  \item[Detection] HSM self-test failure; share key derivation
    refuses to produce a result.
\end{description}

\paragraph{F4 --- Single share material leak.}
\begin{description}[leftmargin=2em,style=nextline,itemsep=0pt]
  \item[Trigger] One operator's share leaves the HSM boundary
    (insider attack, side channel, memory dump).
  \item[Consequence] By T1, no signature forgery yet (leaked share
    counts toward the $t-1$ tolerated minority).
  \item[Recovery] Immediate proactive reshare (\S3.7); the leaked share
    becomes useless one epoch later. The transparency log records the
    rotation event for auditability.
  \item[Detection] HSM audit log shows export attempt; or threat intel
    feed flags the leaked material. This is the hardest mode to
    detect — the operator-side organisational controls (background
    checks, dual-control share access, off-site backup) are the
    primary mitigation.
\end{description}

\paragraph{F5 --- $t-1$ share leak.}
\begin{description}[leftmargin=2em,style=nextline,itemsep=0pt]
  \item[Trigger] $t-1$ operators' shares leak simultaneously (concerted
    insider attack, supply-chain compromise).
  \item[Consequence] Adversary at boundary of T1; needs one more
    share to forge.
  \item[Recovery] Emergency reshare, on-chain key revoke
    (\texttt{Revoke} instruction, \S5.3), and re-register with a fresh
    DKG. All in-flight issuance pauses; existing coins remain valid
    (their signature was produced before the leak).
  \item[Detection] Multiple HSM audit log anomalies in a window;
    correlation with operator personnel events.
\end{description}

\paragraph{F6 --- $t$ or more share leak (catastrophic).}
\begin{description}[leftmargin=2em,style=nextline,itemsep=0pt]
  \item[Trigger] $\geq t$ operators' shares leak (or $t$ operators
    actively collude).
  \item[Consequence] Joint secret $x$ reconstructable via Lagrange
    interpolation (\S3.6); adversary can forge arbitrary coins
    under that joint key. T1 no longer holds for the affected epoch.
  \item[Recovery] On-chain revocation of the compromised keyset
    (\texttt{Revoke} instruction freezes new spends), forensic audit,
    user-side claim against the vault collateral (\S8.6 outlines the
    governance process for slashing the compromised operators' bonds
    if applicable). \emph{No protocol-level rollback}: existing
    transactions in the chain are not reversed.
  \item[Detection] Anomalous issuance volume; on-chain analytics flag
    a spike in newly registered nullifiers without matching
    pre-deposit transfers.
\end{description}

\subsection{User-Side Failure Modes}
\label{sec:failure-user}

\paragraph{F7 --- Lost coin material.}
\begin{description}[leftmargin=2em,style=nextline,itemsep=0pt]
  \item[Trigger] User deletes or loses their wallet's coin record
    (the $(x, Cp, \sigma)$ triple).
  \item[Consequence] The coin's collateral becomes irretrievable
    \emph{unless} the user holds the AEAD-encrypted backup
    (\S6 wallet layer, \texttt{backup.rs}). TARDUS coins are bearer
    instruments: there is no protocol-level recovery from loss.
  \item[Recovery] Restore from \texttt{seal\_backup}/\texttt{open\_backup}
    using the BIP-39 mnemonic-derived master seed (\S6.4 wallet
    backup). If the user has no backup, the coin's lamports remain
    locked in the vault, unspendable, until the protocol's collateral
    expiry rule (operationally: never; the vault retains them
    indefinitely as overcollateralisation).
  \item[Detection] User-side only; the protocol has no view into
    wallet contents.
\end{description}

\paragraph{F8 --- Wallet device theft.}
\begin{description}[leftmargin=2em,style=nextline,itemsep=0pt]
  \item[Trigger] User's device is stolen with unencrypted coin store.
  \item[Consequence] Thief can spend any unspent coin (bearer model).
  \item[Recovery] User submits \texttt{Refresh} transactions for every
    coin they remember (race against the thief). The first to reach
    the on-chain nullifier set wins.
  \item[Detection] User-side only; consider hardware-wallet integration
    (Faz 3 wallet roadmap) for stronger device-level mitigation.
\end{description}

\paragraph{F9 --- Failed refresh session.}
\begin{description}[leftmargin=2em,style=nextline,itemsep=0pt]
  \item[Trigger] Network drop, validator partition (F2), or user
    timeout during the 6-round refresh protocol (\S4).
  \item[Consequence] If the session aborted \emph{before} Round 3
    (mint's challenge), the user retains the old coin intact and may
    retry. If aborted \emph{after} Round 5 (mint's signed reveal),
    the user holds an unfinalised new coin but the old coin is
    nullified on chain — the user must complete Round 6 (off-chain
    unblind) using the persisted refresh-session state from
    \texttt{tardus-client::backup.rs} to recover the new coin.
  \item[Recovery] The wallet persists per-session state on every round
    transition; on restart it reads the state and resumes from the
    last completed round.
  \item[Detection] User-facing wallet timeout; the session state
    encodes its own progress.
\end{description}

\paragraph{F10 --- Refresh session race (network reordering).}
\begin{description}[leftmargin=2em,style=nextline,itemsep=0pt]
  \item[Trigger] Two refresh sessions on the same coin (e.g.\ retry
    of a session whose status is unclear) reach the on-chain program
    in different orders across validators.
  \item[Consequence] First-to-finalise wins; the second aborts with
    \texttt{ERR\_DOUBLE\_SPEND} (T3). User-facing: they hold one valid
    new coin instead of two.
  \item[Recovery] None needed; this is the intended atomic behaviour.
    The CLI's preflight nullifier check (\texttt{tardus devnet refresh})
    surfaces the duplicate before the second TX pays a fee.
  \item[Detection] CLI preflight; or post-finalisation \texttt{Refresh}
    transaction error.
\end{description}

\subsection{Protocol-Level Adversarial Failures}
\label{sec:failure-adversarial}

\paragraph{F11 --- DKG ceremony cheater.}
\begin{description}[leftmargin=2em,style=nextline,itemsep=0pt]
  \item[Trigger] During DKG, one or more operators publish broadcasts
    inconsistent with the shares they send privately (deviating from
    \S3.5).
  \item[Consequence] Honest operators detect the inconsistency via the
    dual-commitment check ($C_{i,0} = A_{i,0}$ in canonical form
    over both bases $G$ and $H$). The ceremony aborts before any
    operator commits to a final share.
  \item[Recovery] Restart DKG with the cheater excluded; publish the
    cheating proof to the transparency log.
  \item[Detection] Automatic during DKG round 2 validation.
\end{description}

\paragraph{F12 --- Reshare cheater (non-zero polynomial).}
\begin{description}[leftmargin=2em,style=nextline,itemsep=0pt]
  \item[Trigger] During proactive reshare (\S3.7), an operator's
    delta-polynomial $\tilde{f}_i$ has $\tilde{f}_i(0) \neq 0$
    (which would shift the joint key).
  \item[Consequence] Honest operators detect via the
    $\tilde{A}_{i,0} = \mathcal{O}$ check (added in v1.4.3 spec
    discovery). The reshare aborts; old shares remain valid for
    the current epoch.
  \item[Recovery] Restart reshare excluding the cheater; operators'
    next-epoch keyset rotation is delayed by one cycle.
  \item[Detection] Automatic during reshare validation.
\end{description}

\paragraph{F13 --- On-chain instruction spam (DoS).}
\begin{description}[leftmargin=2em,style=nextline,itemsep=0pt]
  \item[Trigger] Adversary submits high-volume invalid \texttt{Refresh}
    or \texttt{Withdraw} transactions to exhaust validator compute.
  \item[Consequence] Solana's transaction fee market increases TX cost
    during congestion; honest users pay more but the protocol remains
    available (per-instruction CU is bounded at $\sim$73k, well within
    the 200k default budget, so an attacker cannot cheap-spam).
  \item[Recovery] Solana's fee dynamics auto-rate-limit; no protocol
    action required.
  \item[Detection] Standard Solana RPC metrics.
\end{description}

\paragraph{F14 --- Nullifier-set capacity exhaustion.}
\begin{description}[leftmargin=2em,style=nextline,itemsep=0pt]
  \item[Trigger] The Borsh-encoded \texttt{NullifierSet} PDA reaches its
    8\,KiB allocation cap ($\approx 250$ entries).
  \item[Consequence] New \texttt{Refresh} and \texttt{Withdraw}
    transactions fail with \texttt{ERR\_ACCOUNT\_DATA\_TOO\_SMALL}.
  \item[Recovery] v1.4.13 migration to Light Protocol's compressed
    Merkle tree raises capacity to $2^{32}$ leaves with constant
    on-chain storage (root only). Until then, operators monitor
    nullifier count and reallocate the PDA (a one-time event for
    the lifetime of the deployment, since the BPF Loader Upgradeable
    permits account-data reallocation by the program authority).
  \item[Detection] \texttt{tardus devnet info} reports nullifier count;
    threshold-based alerting at 80\% of capacity.
\end{description}

\subsection{Network-Level Failures}
\label{sec:failure-network}

\paragraph{F15 --- Solana network outage.}
\begin{description}[leftmargin=2em,style=nextline,itemsep=0pt]
  \item[Trigger] Solana mainnet halts (block production stops).
  \item[Consequence] No new \texttt{Refresh} or \texttt{Withdraw}
    transactions land; existing coins remain valid; the mint
    committee continues to issue off-chain (issuance does not
    require on-chain TX). Withdrawals stall.
  \item[Recovery] Resume when Solana production resumes; pending
    transactions are picked up automatically by the validator
    daemon's retry queue.
  \item[Detection] Solana RPC reports stale slot height; transparency
    log records the pause.
\end{description}

\paragraph{F16 --- Solana chain reorganisation past finality.}
\begin{description}[leftmargin=2em,style=nextline,itemsep=0pt]
  \item[Trigger] An unprecedented Solana consensus event reorganises
    blocks at depth $\geq 32$ (the safety boundary in A6).
  \item[Consequence] Catastrophic for any Solana program; TARDUS
    inherits the network's safety property. T3 and T6 fail if A6
    fails.
  \item[Recovery] Coordinated with Solana ecosystem; nothing TARDUS
    can do unilaterally.
  \item[Detection] Solana RPC reports finalised-block divergence;
    network-wide alert.
\end{description}

\paragraph{F17 --- Token-2022 program bug (post-v1.4.12).}
\begin{description}[leftmargin=2em,style=nextline,itemsep=0pt]
  \item[Trigger] A bug in Solana's Token-2022 program affects the
    Confidential Mint extension (v1.4.12+).
  \item[Consequence] Vault SPL operations may revert or, worse, leak
    amounts that should be encrypted. v1.4.12 will document the
    explicit dependence on Token-2022 invariants.
  \item[Recovery] Coordinate with Solana Labs; pause vault operations
    via \texttt{Revoke} until the upstream fix is deployed.
  \item[Detection] Standard SPL ecosystem alerting.
\end{description}

\paragraph{F18 --- Malicious RPC node.}
\begin{description}[leftmargin=2em,style=nextline,itemsep=0pt]
  \item[Trigger] User's RPC provider returns stale or fabricated
    account data.
  \item[Consequence] CLI preflight check (\texttt{query-nullifier})
    could miss an already-inserted nullifier, leading to a wasted
    TX fee when the on-chain check catches it. The on-chain check
    always wins (T3), so no double-spend risk.
  \item[Recovery] Multi-RPC verification (Faz 3 wallet roadmap); or
    user-side fallback to the validator committee's own RPC.
  \item[Detection] Cross-RPC discrepancy alarms; user-facing warning.
\end{description}

\subsection{Relay-Layer Failure Modes}
\label{sec:failure-relay}

The relay layer (\Cref{sec:relay}) introduces its own operational
F-modes catalogued as R1--R7 in \Cref{sec:relay-failure-modes}.
Cross-referenced summary:

\paragraph{R1 --- Relay process crash.} F1-class extension: with
v5.4 \texttt{SQLite} persistence (\Cref{sec:relay-storage}) the
in-flight inbox survives restart; the validator-side equivalent of
F1's ``operator restart'' applies directly.

\paragraph{R2 --- Relay disk full.} F1-class extension: \texttt{SQLite}
writes fail $\to$ POST returns HTTP 5xx $\to$ sender's wallet
retries with backoff. No coin-side state damage.

\paragraph{R3 --- Relay operator payload inspection.} Mitigated by
the v5.6 sealed-box payload format (\Cref{sec:relay-payload}):
the relay cannot recover coin material under T7. A relay-side audit
still sees recipient pubkey, timestamps, ciphertext sizes (traffic
metadata) — Tor/I2P circuiting is the deployment-time defence.

\paragraph{R4 --- Relay log seizure.} Same mitigation as R3: logs
hold only opaque ciphertext + metadata.

\paragraph{R5 --- Inbox exhaustion attack.} F13-class extension: a
malicious sender deposits up to \texttt{max\_per\_recipient} entries
on the victim's inbox. The v5.6 mitigation is the per-recipient cap
(503 on overflow) + TTL-based eviction. Long-term mitigation is a
fee market (deferred to v5.7).

\paragraph{R6 --- Recipient secret leaked.} F8-class extension to
the wallet section: leak of the recipient's \texttt{recv\_sk}
(\Cref{sec:wallet-recv-identity}) lets the attacker decrypt all
ciphertexts addressed to that pubkey, past and future. The wallet
mnemonic is therefore the security root for both coin custody and
P2P payload privacy.

\paragraph{R7 --- Sender IP fingerprinting.} Not modelled in v5; a
relay-side packet sniffer can link a recipient pubkey to a sender IP.
Production deployments are expected to circuit sender wallets
through Tor or a privacy-respecting VPN before reaching the relay.

\subsection{HSM-Backend Failure Modes}
\label{sec:failure-hsm}

The v2.11 PKCS\#11 backend
(\Cref{sec:mint-validator-daemon}) introduces failure modes that the
v2.1 file backend did not have. These extend F3 (HSM failure):

\paragraph{F3a --- PKCS\#11 module path missing.}
\begin{description}[leftmargin=2em,style=nextline,itemsep=0pt]
  \item[Trigger] \texttt{--hsm-module} points to a missing or
    incompatible \texttt{.so}.
  \item[Consequence] Validator boot fails with explicit
    \texttt{Error::Config}; no shares ever loaded.
  \item[Recovery] Fix path; rotate operator if the wrong module was
    deployed.
\end{description}

\paragraph{F3b --- PKCS\#11 user PIN rotated upstream.}
\begin{description}[leftmargin=2em,style=nextline,itemsep=0pt]
  \item[Trigger] HSM operator rotates the user PIN through the
    HSM-vendor management console without updating the validator's
    config.
  \item[Consequence] Validator boot fails at the login step
    (\texttt{C\_Login} returns \texttt{CKR\_PIN\_INCORRECT}); no
    silent fallback.
  \item[Recovery] Sync new PIN; validator restart.
\end{description}

\paragraph{F3c --- Wrap key deleted on the HSM.}
\begin{description}[leftmargin=2em,style=nextline,itemsep=0pt]
  \item[Trigger] Operator misclick or HSM tenancy event removes the
    AES wrap key (\texttt{CKA\_LABEL = tardus-wrap-v1}).
  \item[Consequence] All on-disk \texttt{share\_*.bin} files become
    permanently undecryptable; the affected validator must
    re-participate in a reshare to obtain a fresh share. This is
    F4-class for the affected validator but does not compromise the
    threshold.
  \item[Recovery] Reshare under existing \texttt{joint\_pk}; the
    re-keyed validator picks up its new share from the reshare
    transcript.
\end{description}

\paragraph{F3d --- HSM clock or session timeout.}
\begin{description}[leftmargin=2em,style=nextline,itemsep=0pt]
  \item[Trigger] HSM enforces session lifetime; the
    \texttt{Pkcs11ShareStore}'s long-held session is invalidated.
  \item[Consequence] Next \texttt{C\_Encrypt}/\texttt{C\_Decrypt}
    fails with \texttt{CKR\_SESSION\_HANDLE\_INVALID}.
  \item[Recovery] v2.12 will add automatic session re-open on
    timeout; v2.11 requires validator restart.
\end{description}

\subsection{Detection and Monitoring Discipline}
\label{sec:failure-monitoring}

The full failure-mode catalogue above is operationalised through three
observation channels:

\begin{enumerate}[leftmargin=2em,itemsep=0pt]
  \item \textbf{Transparency log.} Append-only, hash-chained log
    published by the mint committee, recording every DKG round,
    every reshare event, every keyset registration/revocation, and
    every operator heartbeat. Anyone can replay the log to verify
    committee behaviour. F1, F2, F4, F5, F11, F12 surface here.
  \item \textbf{On-chain analytics.} Standard Solana indexers query
    the program's account states. \texttt{tardus devnet info}
    (\S6.5) is the minimal CLI front for this; richer dashboards
    (Faz 3) will track nullifier-rate, vault-balance, and
    keyset-registry growth. F6, F13, F14, F18 surface here.
  \item \textbf{Operator alerting.} Per-validator daemon (Faz 2)
    Prometheus-style metrics on heartbeat, share count, ceremony
    participation, HSM health. F1, F2, F3, F4 surface here first.
\end{enumerate}

\subsection{Governance and Open Problems}
\label{sec:failure-open}

\paragraph{Slashing.} The current v1 protocol does not implement
operator-bond slashing; F4-F6 recovery depends on out-of-band
governance (e.g.\ legal action against an operator who leaked
material). A v2 extension under consideration: operators post a
collateral bond into a separate Solana PDA; the
\texttt{Revoke} instruction can be augmented with a slashing
parameter that diverts the bond to the vault as compensation. This is
deferred to Faz 5.

\paragraph{Cross-chain bridge attacks.} Coins denominated in
bridged assets (e.g.\ Wormhole-wrapped tokens) inherit the bridge's
security posture. TARDUS itself cannot defend against a bridge
compromise. This will be addressed in the Faz 5 mainnet asset-class
policy: only assets with native Solana issuance (SOL, Token-2022 mints
with the issuer's signature attestation) qualify for the canonical
TARDUS keyset; bridged assets get separate keysets with explicit
risk disclosure.

\paragraph{Post-quantum migration.} A1 (ECDLP) is broken by a
sufficiently large quantum computer (Shor's algorithm, $\sim 4000$
logical qubits). TARDUS will need a post-quantum signature primitive
in that era; current candidates (Dilithium, Falcon) have $\sim 2$\,KB
signatures, incompatible with the per-coin storage budget. The
mitigation will be a parallel post-quantum keyset (separate
\texttt{KeysetEntry} with a post-quantum algorithm marker), with the
classical keyset deprecated over a multi-year migration window. See
\S5 for the keyset registry's extensibility.

\paragraph{Async secure messaging.} Faz 2's validator daemon needs
authenticated, replay-protected, end-to-end-encrypted messaging
between operators. The current design assumes TLS+mTLS over a
private network; a fully decentralised alternative
(libp2p, Iroh) is under evaluation.

\subsection{Summary Table}
\label{sec:failure-summary}

\begin{center}
\small
\begin{tabular}{@{}llll@{}}
\toprule
ID & Mode & Severity & Recovery class \\
\midrule
F1  & Single validator crash       & Low      & Operator restart (v2.X)         \\
F2  & Validator network partition  & Medium   & Auto-quorum-pause; manual recon \\
F3  & HSM failure                  & Low      & Reshare + replace               \\
F4  & Single share leak            & Medium   & Proactive reshare               \\
F5  & $t-1$ share leak             & High     & Emergency reshare + revoke      \\
F6  & $\geq t$ share leak          & Catastrophic & On-chain revoke + audit     \\
F7  & Lost coin material           & High (user) & AEAD backup restore          \\
F8  & Wallet theft                 & High (user) & Race-refresh; HW wallet      \\
F9  & Refresh session abort        & Low      & Wallet auto-resume              \\
F10 & Refresh race                 & None     & First-finalised wins (T3)       \\
F11 & DKG ceremony cheater         & Low      & Auto-detect + restart           \\
F12 & Reshare cheater              & Low      & Auto-detect + restart           \\
F13 & Instruction spam (DoS)       & Low      & Solana fee market               \\
F14 & Nullifier capacity exhausted & Medium   & Reallocation; v1.4.13 Light tree \\
F15 & Solana network outage        & Low      & Wait + auto-retry               \\
F16 & Solana finality break        & Catastrophic & Ecosystem-wide              \\
F17 & Token-2022 bug (v1.4.12+)    & Variable & Pause + upstream fix            \\
F18 & Malicious RPC                & Low      & Multi-RPC verification          \\
\bottomrule
\end{tabular}
\end{center}

Of the 18 catalogued modes, 14 have automated detection and protocol-level
recovery (F1-F3, F9-F18); 4 require human-in-the-loop response (F4-F8).
This is the explicit operational surface the validator daemon (Faz 2)
must instrument.

\section{Relay Layer}
\label{sec:relay}

This section specifies TARDUS's optional relay layer (\Cref{sec:wallet} §6
referenced this layer as the canonical \emph{user-to-user delivery
mechanism} for off-chain coin transmission). The relay is a stateless,
TTL-bound, recipient-keyed mailbox: senders deposit ciphertexts addressed
to a recipient public key, and recipients poll. The relay never holds
key material and (under the v5.6 sealed-box payload format) never sees
plaintext coin material.

\subsection{Design Objectives}
\label{sec:relay-objectives}

\begin{description}[leftmargin=2em,style=nextline,itemsep=0pt]
  \item[O1 (recipient confidentiality)] A relay operator must not learn
    coin material (secrets, signatures, denomination) even on a full
    payload audit, except via decryption of the sealed-box payload —
    which requires the recipient's ed25519 secret key.
  \item[O2 (sender anonymity)] The relay must not authenticate the
    sender. The POST endpoint accepts unauthenticated requests; the
    only network-level identifier the relay sees is the connecting IP
    (mitigated at deployment by Tor/I2P or VPN tunnelling — out of
    scope for v5).
  \item[O3 (bounded resources)] Per-recipient and per-payload caps
    bound the relay's memory and disk usage; an oversized payload or
    a full inbox is rejected at the HTTP layer with 400 / 503.
  \item[O4 (eventual delivery)] Messages have a TTL (default 7 days);
    expired entries are pruned by a background task. The TTL forces
    recipients to poll within a bounded window; longer-term storage is
    the wallet's responsibility.
  \item[O5 (operator agnosticism)] The relay never inspects the
    payload's contents and never participates in the mint or refresh
    protocols. A relay outage delays delivery but cannot forge,
    double-spend, or otherwise compromise coin security.
\end{description}

\subsection{Wire Protocol}
\label{sec:relay-protocol}

The relay exposes three HTTP endpoints, all under the canonical path
\verb|/inbox/{recipient_pk_hex}|.

\paragraph{POST /inbox/\{recipient\_pk\_hex\}} — anonymous deposit.

\begin{description}[leftmargin=2em,itemsep=0pt]
  \item[Request]
    \verb|{ "payload_hex": "<bytes>", "ttl_secs": optional u64 }|
  \item[Response 200]
    \verb|{ "id": "<16-byte hex>", "payload_hex": "...", "received_at_unix_ms": u64, "expires_at_unix_ms": u64 }|
  \item[Response 400] payload exceeds \verb|--max-payload-bytes|, or
    \verb|recipient_pk_hex| not 32 bytes of hex.
  \item[Response 503] inbox full for this recipient
    (\verb|--max-per-recipient| reached).
\end{description}

\paragraph{GET /inbox/\{recipient\_pk\_hex\}} — recipient polls.

\begin{description}[leftmargin=2em,itemsep=0pt]
  \item[Response 200]
    \verb|{ "messages": [{ id, payload_hex, received_at_unix_ms, expires_at_unix_ms }, ...] }|
  \item[Response 400] \verb|recipient_pk_hex| malformed.
\end{description}

\paragraph{DELETE /inbox/\{recipient\_pk\_hex\}/\{message\_id\}} —
recipient marks consumed.

\begin{description}[leftmargin=2em,itemsep=0pt]
  \item[Response 200] \verb|{ "removed": bool }|
\end{description}

\paragraph{Operator endpoints.} \verb|GET /health| (uptime, deposits
total, fetches total, messages inflight) and \verb|GET /info|
(operator name, bind addr, caps, default TTL) for monitoring; no
mutation, no auth required.

\subsection{Payload Format}
\label{sec:relay-payload}

The wire-level \verb|payload_hex| is opaque to the relay. Two formats
coexist in v5.6:

\paragraph{Plain JSON (v5.3, deprecated for production).} The wallet
JSON-encodes
\[
  \{\,\texttt{coin\_secret},\,\texttt{coin\_pubkey},\,\texttt{coin\_signature},\,\texttt{denom},\,\texttt{memo}\,\}
\]
and hex-encodes the result. A relay operator with read access to the
inbox can recover all coin material. Used only when the recipient is
unable to consume sealed-box payloads (e.g. lightweight read-only
audit clients).

\paragraph{Sealed box (v5.6, default for production).} The payload is
\[
  \texttt{sealed\_blob} \;=\; \texttt{epk}\;(32)\;\|\;\texttt{ChaCha20\text{-}Poly1305(plaintext)}
\]
where the plaintext is the same JSON as above, the AEAD key is derived
via the NaCl \verb|crypto_box_seal| construction adapted for ed25519
recipient keys:

\begin{enumerate}[leftmargin=2em,itemsep=0pt]
  \item Sender samples ephemeral X25519 keypair $(esk, epk)$.
  \item Recipient's ed25519 pubkey $\to$ Montgomery form via RFC 7748
        §5 (the Edwards-to-Montgomery birational map).
  \item Shared secret $ss = esk \cdot \mathsf{recipient}_{X25519}$.
  \item $(K, N) =
        \mathsf{HKDF\text{-}SHA\text{-}256}(\,\mathsf{salt}=\texttt{"TARDUS-sealed-box-v1"},\,
                                            \mathsf{ikm}=ss,\,
                                            \mathsf{info}=epk \,\|\, \mathsf{recipient}_{X25519}\,)$.
  \item $\mathsf{ciphertext} =
        \mathsf{ChaCha20\text{-}Poly1305}(K, N, \mathsf{plaintext})$.
\end{enumerate}

Recipient inverts steps 1–4 using its own ed25519 secret key (also
canonical scalar, converted to X25519 scalar) and then runs the
matching $\mathsf{ChaCha20\text{-}Poly1305\text{-}Decrypt}$.

\subsection{Storage Backends}
\label{sec:relay-storage}

The relay's \verb|InboxStore| dispatches to one of two backends, chosen
at boot via the \verb|--storage-path| CLI flag.

\paragraph{In-memory backend (default).} \verb|HashMap<RecipientPk, Vec<StoredMessage>>|
behind a \verb|std::sync::Mutex|. Messages are lost on restart. Suited
to ephemeral testing and short-lived federation peers.

\paragraph{SQLite backend (production).} Bundled SQLite (no system
dependency), schema:

\begin{verbatim}
CREATE TABLE messages (
    id                  TEXT PRIMARY KEY,
    recipient           BLOB NOT NULL,
    payload_hex         TEXT NOT NULL,
    received_at_unix_ms INTEGER NOT NULL,
    expires_at_unix_ms  INTEGER NOT NULL
);
CREATE INDEX idx_recipient ON messages(recipient);
CREATE INDEX idx_expires   ON messages(expires_at_unix_ms);
\end{verbatim}

The two indices give $O(\log n)$ recipient-bucket scans and $O(\log n)$
TTL-prune scans even at $10^6+$ inflight messages.

\paragraph{TTL pruning.} A Tokio background task runs every $60\,\text{s}$
and issues
\verb|DELETE FROM messages WHERE expires_at_unix_ms <= now|
on SQLite or retains-by-deadline on the in-memory backend. Prune is
side-effect-free: no event log emission, no recipient notification.

\subsection{Transport Security}
\label{sec:relay-tls}

Production relays MUST run with TLS enabled. The daemon supports HTTPS
via \verb|axum-server| + \verb|rustls 0.23| (ring backend); operator
provides \verb|--tls-cert| and \verb|--tls-key| at boot. A relay with
neither flag set serves plain HTTP and refuses to accept either flag in
isolation (fail-fast misconfiguration check). v5.7 will add mTLS
support for recipient cert-based authentication on GET / DELETE.

\subsection{Security Argument}
\label{sec:relay-security}

\begin{theorem}[Sealed-box payload confidentiality]
\label{thm:relay-confidentiality}
Under assumption A1 (ECDLP) and the IND-CCA security of
ChaCha20-Poly1305, an adversary controlling the relay operator and
observing all inbox contents learns no information about the plaintext
coin material beyond the bit-length of each ciphertext, except with
probability $\leq \mathsf{Adv}^{\mathsf{IND\text{-}CCA}}_{\mathsf{ChaCha20\text{-}Poly1305}}(\lambda)
+ \mathsf{Adv}^{\mathsf{ECDLP}}_{X25519}(\lambda) + \negl(\lambda)$.
\end{theorem}

\paragraph{Reduction sketch.} The reduction is the standard
\verb|crypto_box_seal| argument: the X25519 ECDH shared secret is
indistinguishable from random under DDH (which reduces to ECDLP for
the Curve25519 family), HKDF-SHA-256 is a PRF under ROM, and
ChaCha20-Poly1305 is IND-CCA secure. Composition is via standard
hybrid arguments.

\paragraph{What the relay still learns.} The relay sees, for each
deposit: the recipient pubkey (in the URL path), the deposit timestamp,
the message TTL, and the ciphertext length. Traffic-analysis resistance
against these side channels is a deployment concern (Tor/I2P circuit,
fixed-size payload padding) and out of scope for v5.

\paragraph{What an inbox-leak does \emph{not} compromise.}
Critically, a relay compromise does not give the adversary:
(a) the recipient's ed25519 secret key (held only by the recipient's
wallet), (b) the sender's identity (no auth on POST), (c) the validator
committee's joint secret (off-chain, held by validator daemons), or
(d) any prior payments (no logging of forwarded message contents).

\subsection{Failure Mode Mapping}
\label{sec:relay-failure-modes}

Cross-reference with \Cref{sec:failure-modes}: the relay layer's
operational failure modes are extensions of the wallet/operator F-modes,
not new categories.

\begin{center}
\small
\begin{tabular}{@{}lll@{}}
\toprule
ID & Relay-layer scenario & Existing F-mode mapping \\
\midrule
R1 & Relay process crash       & F1-class (restart with persistent DB) \\
R2 & Relay disk full           & F1-class extended (SQLite write fails → 5xx) \\
R3 & Relay operator inspection & R-3 below (sealed-box mitigation) \\
R4 & Relay log seizure         & R-3 same                                \\
R5 & Inbox exhaustion attack   & F13-class (per-recipient cap, fee market) \\
R6 & Recipient secret leaked   & F8-class (wallet theft) extended         \\
R7 & Sender IP fingerprinting  & not modelled in v5 (Tor circuit at deploy) \\
\bottomrule
\end{tabular}
\end{center}

\paragraph{R-3 — relay operator payload inspection.} Mitigated by
the v5.6 sealed-box payload format. A compromised or hostile relay
operator can still observe recipient pubkeys and message metadata
(deposit timing, ciphertext size, TTL choice), but cannot recover
any coin secret or signature.

\subsection{Federation (Forward Notice)}
\label{sec:relay-federation}

A single relay is a single point of failure for delivery. The v5.7
federation extension will define a gossip protocol for relay-to-relay
forwarding so a recipient can poll any peer relay and recover messages
deposited at any federated relay. This is intentionally deferred until
the single-relay deployment pattern is stable in production.


% --- appendix ---
\appendix

\section{Reference Test Vectors}
\label{sec:appendix-vectors}

\paragraph{Purpose.} This appendix lists deterministic test
vectors used by the cross-language (Rust + TypeScript)
reference implementations. Any production or audit
implementation MUST reproduce these outputs bit-equal from the
given inputs. Source: \texttt{crates/tardus-wallet/src/mnemonic.rs}
+ \texttt{ts-sdk/src/mnemonic-rust-compat.ts}.

License: TARDUS-PROPRIETARY-1.0.

\subsection{BIP-39 \texorpdfstring{$\to$}{->} master seed}
\label{sec:vec-master-seed}

The canonical BIP-39 test mnemonic from the original BIP-39
specification (Trezor reference, 12 words):

\begin{center}\small
\begin{tabular}{rl}
\toprule
\texttt{mnemonic}      & \texttt{abandon abandon abandon abandon abandon abandon} \\
                       & \texttt{abandon abandon abandon abandon abandon about}   \\
\texttt{passphrase}    & (empty string)                                            \\
\texttt{master\_seed (TARDUS, first 32 bytes of PBKDF2 seed)} & \\
                       & \texttt{5eb00bbddcf069084889a8ab9155568165f5c453ccb85e70811aaed6f6da5fc1} \\
\bottomrule
\end{tabular}
\end{center}

Reproduce with:
\begin{itemize}[leftmargin=2em,itemsep=0pt]
  \item Rust: \texttt{tardus\_wallet::derive\_master\_seed(\&mnemonic, "")}.
  \item TS: \texttt{deriveMasterSeed(phrase, '')} (v0.1, X25519 path) OR
        \texttt{deriveMasterSeedRustCompat(phrase, '')} (v0.2,
        Rust-compat).
\end{itemize}

\subsection{Master seed \texorpdfstring{$\to$}{->} receiving keypair (Rust + TS v0.2)}
\label{sec:vec-recv-keypair}

\begin{center}\small
\begin{tabular}{rl}
\toprule
\texttt{master\_seed (input)} & \texttt{5eb00bbddcf069084889a8ab9155568165f5c453ccb85e70811aaed6f6da5fc1} \\
\texttt{HKDF salt}            & ASCII \texttt{"TARDUS-recv-id-v1"} \\
\texttt{HKDF info}            & ASCII \texttt{"ed25519-keypair"} \\
\texttt{HKDF output length}   & 64 bytes (HKDF-SHA-512) \\
\texttt{sk (canonical Ed25519 scalar, mod $\ell$)} & \\
                              & \multicolumn{1}{l}{\textit{Deterministic; bit-equal Rust} } \\
                              & \multicolumn{1}{l}{\textit{$\to$ TS v0.2.}} \\
\texttt{pk (32-byte Edwards-compressed point)} & \\
                              & \multicolumn{1}{l}{\textit{$= (sk \cdot G_{\mathrm{ed25519}}).\mathsf{compress}()$.}} \\
\bottomrule
\end{tabular}
\end{center}

Cross-validate by running \texttt{cargo test -p tardus-wallet
test\_canonical\_master\_seed} (Rust) and
\texttt{node --test --experimental-strip-types ts-sdk/test/mnemonic.test.ts}
(TS v0.1) or
\texttt{ts-sdk/test/rust-compat.test.ts} (TS v0.2). The first 20
hex characters of \texttt{master\_seed} are asserted in both
languages' test suites as \texttt{5eb00bbddcf069084889}.

\subsection{Sealed-box AEAD wire format}
\label{sec:vec-sealed-box-wire}

The wire format (32-byte ephemeral public key followed by
AEAD ciphertext + 16-byte Poly1305 tag) is shared between Rust
and TS v0.2.

\begin{center}\small
\begin{tabular}{rl}
\toprule
\texttt{offset 0..32}     & ephemeral public key $epk$ \\
                          & (Curve25519 u-coordinate, little-endian) \\
\texttt{offset 32..(N+48)} & \texttt{ChaCha20-Poly1305}$_{K,N}$\texttt{(plaintext)} \\
                          & where $K$ (32 B) and $N$ (12 B) are HKDF-derived. \\
\bottomrule
\end{tabular}
\end{center}

HKDF inputs (both Rust and TS v0.2 MUST agree byte-for-byte):

\begin{center}\small
\begin{tabular}{rl}
\toprule
\texttt{salt}             & ASCII \texttt{"TARDUS-sealed-box-v1"} (20 bytes) \\
\texttt{ikm (input keying material)} & ECDH shared secret \\
                          & ($= esk \cdot pk_R^{(\mathrm{X25519})}$, 32 bytes) \\
\texttt{info for key}     & ASCII \texttt{"chacha20poly1305-key"} $\,\|\,$ $epk$ $\,\|\,$ $pk_R^{(\mathrm{X25519})}$ \\
\texttt{info for nonce}   & ASCII \texttt{"chacha20poly1305-nonce"} $\,\|\,$ $epk$ $\,\|\,$ $pk_R^{(\mathrm{X25519})}$ \\
\texttt{key output length} & 32 bytes \\
\texttt{nonce output length} & 12 bytes \\
\bottomrule
\end{tabular}
\end{center}

\textbf{TS v0.1 caveat.} The TS v0.1 path (\texttt{src/sealed-box.ts})
uses X25519-clamped scalar derivation per RFC 7748 §5; its
on-the-wire ciphertexts are NOT decryptable by the Rust
reference. The v0.2 path
(\texttt{src/sealed-box-rust-compat.ts}) uses the unclamped
Curve25519 Montgomery ladder
(\texttt{src/montgomery.ts}) and IS wire-byte-compatible.

\subsection{Nullifier computation}
\label{sec:vec-nullifier}

The on-chain nullifier of a coin with public commitment $Cp$:

\[
  \mathsf{nullifier} \;=\; \mathsf{SHA\text{-}256}(\mathtt{"TARDUS-nullifier-v1"} \,\|\, Cp).
\]

Reproduce with
\texttt{tardus\_program::processor::compute\_nullifier(\&cp\_bytes)}.

\subsection{PDA seeds}
\label{sec:vec-pda-seeds}

Solana program-derived addresses for TARDUS PDAs are computed
via \texttt{Pubkey::find\_program\_address(\&[\&str, ...], \&program\_id)}.
Seeds:

\begin{center}\small
\begin{tabular}{rl}
\toprule
\texttt{KeysetRegistry}  & \texttt{["tardus", "keyset-registry"]}                       \\
\texttt{NullifierTree}   & \texttt{["tardus", "nullifier-tree"]}                        \\
\texttt{Vault}           & \texttt{["tardus", "vault", \&denom.to\_le\_bytes()]}          \\
\texttt{SponsorPool}     & \texttt{["tardus", "sponsor-pool"]}                          \\
\bottomrule
\end{tabular}
\end{center}

On devnet with the canonical program ID
\texttt{AmY1ysgQyCC6CmorXkrNkogBSHmomy4kbsPziAYUx47u}, the
KeysetRegistry PDA resolves to
\texttt{4XJ8fmJgJEc9QhNWgNdw1xJHzs3cEZ6JhM746KWDpHK2}.
Reproduce with
\texttt{tardus devnet info --rpc https://api.devnet.solana.com}.

\subsection{Ed25519 precompile envelope}
\label{sec:vec-ed25519-precompile}

The TARDUS \texttt{Refresh} and \texttt{Withdraw} instructions
require an immediately-preceding ed25519 precompile instruction
with the following 144-byte data layout (single-signature
variant):

\begin{center}\small
\begin{tabular}{rl}
\toprule
\texttt{offset 0..1}    & \texttt{num\_signatures = 1}      \\
\texttt{offset 1..2}    & \texttt{padding = 0}              \\
\texttt{offset 2..4}    & \texttt{sig\_offset = 16} (u16 LE) \\
\texttt{offset 4..6}    & \texttt{sig\_ix\_idx = 0xFFFF} (u16 LE) \\
\texttt{offset 6..8}    & \texttt{pk\_offset = 80}          \\
\texttt{offset 8..10}   & \texttt{pk\_ix\_idx = 0xFFFF}     \\
\texttt{offset 10..12}  & \texttt{msg\_offset = 112}        \\
\texttt{offset 12..14}  & \texttt{msg\_size = 32}           \\
\texttt{offset 14..16}  & \texttt{msg\_ix\_idx = 0xFFFF}    \\
\texttt{offset 16..80}  & 64-byte signature $\sigma$ on coin $Cp$ \\
\texttt{offset 80..112} & 32-byte mint joint pubkey $\mathsf{joint\_pk}$ \\
\texttt{offset 112..144} & 32-byte coin commitment $Cp$ (the message) \\
\bottomrule
\end{tabular}
\end{center}

Reproduce with
\texttt{crates/tardus-cli/src/commands/devnet.rs::build\_ed25519\_precompile\_data}
or
\texttt{crates/tardus-wallet-gui/src/runtime.rs::withdraw\_on\_devnet}.

\subsection{\texorpdfstring{Solana realloc 10\,KB per-instruction cap (discovery \#10)}{Solana realloc 10 KB per-instruction cap (discovery 10)}}
\label{sec:vec-realloc-cap}

\texttt{AccountInfo::realloc(new\_size, true)} is bounded by
\texttt{MAX\_PERMITTED\_DATA\_INCREASE = 10240 bytes} per single
instruction call. Growing a PDA from 1024 bytes to 65536 bytes
therefore requires
$\lceil (65536 - 1024) / 10240 \rceil = 7$ sequential
\texttt{ResizeAccount} instructions, each within its own
transaction (or chained in a single transaction with multiple
program calls). The CLI's
\texttt{tardus devnet resize-nullifier-tree} subcommand handles
the multi-transaction batching transparently; operators see a
single high-level target size.


% --- bibliography ---
\newpage
\bibliography{refs}

\end{document}
